%% (c) 2026 Brayden Sanders / 7Site LLC, Arkansas. All rights reserved.
%% For human use only. No commercial or government use without written agreement from 7Site LLC.

\documentclass[11pt]{article}
\usepackage{amsmath,amssymb,amsthm,mathtools}
\usepackage[margin=1in]{geometry}
\usepackage{hyperref}

\newtheorem{theorem}{Theorem}[section]
\newtheorem{lemma}[theorem]{Lemma}
\newtheorem{proposition}[theorem]{Proposition}
\newtheorem{definition}[theorem]{Definition}
\newtheorem{corollary}[theorem]{Corollary}
\newtheorem{remark}[theorem]{Remark}
\newtheorem{hypothesis}[theorem]{Hypothesis}
\newtheorem{conjecture}[theorem]{Conjecture}

\DeclareMathOperator{\rk}{rk}
\DeclareMathOperator{\ord}{ord}
\DeclareMathOperator{\Reg}{Reg}
\DeclareMathOperator{\Sha}{Sha}
\DeclareMathOperator{\Sel}{Sel}

\title{Elliptic Curve Coherence:\\
TIG Alignment of Analytic Rank and Algebraic Rank\\[6pt]
\large Coherence Field v1.0 --- Problem P5 (BSD)}
\author{Brayden Sanders / 7Site LLC}
\date{February 2026}

\begin{document}
\maketitle

% ==============================================================================
\begin{abstract}
% ==============================================================================

We reformulate the Birch and Swinnerton-Dyer conjecture as a coherence
measurement problem within the Coherence Field framework.  The TIG
operator path $T_1 \to T_2 \to T_5 \to T_7 \to T_9$ (structure $\to$
dual $\to$ feedback $\to$ alignment $\to$ completion) maps the BSD conjecture
onto a defect functional $\delta_{\mathrm{BSD}}(E)$ that measures the mismatch
between the analytic side---the order of vanishing of $L(E,s)$ at $s=1$ and its
leading Taylor coefficient---and the arithmetic side---the Mordell--Weil rank,
Tate--Shafarevich group, regulator, and local Tamagawa data.  We prove that
$\delta_{\mathrm{BSD}}(E) = 0$ if and only if the full BSD conjecture holds for
$E$ (Lemma MC-BSD), conditional on finiteness of $\Sha(E)$ and Euler system
availability.  For analytic rank $r_{\mathrm{an}} \leq 1$, the hypotheses are
satisfied by the theorems of Gross--Zagier, Kolyvagin, and Kato, and
$\delta_{\mathrm{BSD}} = 0$ is verified unconditionally.  For
$r_{\mathrm{an}} \geq 2$, we identify two critical gaps---the absence of
higher-rank Euler systems and the open finiteness of $\Sha$---and characterize
the locked frontier defect $\delta = 1.3$ measured by CK at depth L24.

CK measurements confirm: calibration curves (rank-matched, $y^2 = x^3 - x$)
produce $\delta = 0.0$ exactly; frontier probes (rank-mismatched) produce
$\delta = 1.3$, locked across all seeds with zero coefficient of variation.

\medskip
\noindent\textit{CK measures.  CK does not prove.}
\end{abstract}

% ==============================================================================
\section{Introduction}
\label{sec:intro}
% ==============================================================================

This paper is part of the Coherence Field v1.0 (February 2026), a
systematic program that reformulates each Clay Millennium Problem as a
coherence measurement within the TIG (Thesis--Integration--Genesis) operator
algebra and the SDV (Structure--Defect--Verification) axiom system.

The Birch and Swinnerton-Dyer conjecture occupies a distinguished position
among the Millennium Problems: it is the only one for which substantial partial
results exist.  The theorems of Gross--Zagier \cite{GZ} and Kolyvagin
\cite{Koly} establish BSD for elliptic curves $E/\mathbb{Q}$ with analytic rank
$r_{\mathrm{an}} \leq 1$, and the work of Kato \cite{Kato}, Skinner--Urban
\cite{SU}, and Bhargava--Shankar \cite{BS} provides further arithmetic and
statistical confirmation.  The conjecture remains fully open for
$r_{\mathrm{an}} \geq 2$.

Our contribution is threefold.

\begin{enumerate}
  \item \textbf{Coherence reformulation.}  We define the BSD coherence defect
    $\delta_{\mathrm{BSD}}(E)$, a non-negative real number that vanishes
    precisely when BSD holds for $E$.  The defect decomposes into a rank term
    $|r_{\mathrm{an}} - r|$ and a leading-coefficient term measuring the gap
    between the $L$-function Taylor coefficient and the BSD arithmetic constant.

  \item \textbf{TIG operator path.}  We map the logical structure of BSD onto
    the TIG path $T_1 \to T_2 \to T_5 \to T_7 \to T_9$, identifying the
    feedback operator $T_5$ as the critical link between analytic and arithmetic
    data.  Each operator in the path corresponds to a step in the proof
    skeleton: lattice structure ($T_1$), Pontryagin duality ($T_2$),
    analytic--arithmetic mismatch ($T_5$), rank alignment ($T_7$), and terminal
    coherence ($T_9$).

  \item \textbf{Gap identification.}  We give a precise accounting of which
    steps in the proof skeleton are complete, partially complete, or open, with
    explicit labels \textbf{STANDARD}, \textbf{PARTIALLY COMPLETE}, and
    \textbf{CRITICAL GAP}.  The two critical gaps for $r_{\mathrm{an}} \geq 2$
    are: (i)~the absence of Euler systems that bound Selmer groups, and
    (ii)~the open finiteness of $\Sha(E)$.
\end{enumerate}

\noindent
\textbf{Delta signature}: \texttt{4b5637bfdcd09a00}.

\noindent
\textbf{Problem class}: Affirmative ($\delta \to 0$ when ranks match).

% ── Invention vs Invariant (v1.5) ──
\subsection*{Methodological Note: Invention vs.\ Invariance}
\addcontentsline{toc}{subsection}{Methodological Note: Invention vs.\ Invariance}

The BSD conjecture is entangled with layers of \emph{free invention}:
specific $L$-series representations ($L(E, s)$ via Euler product, modular form,
or Galois representation), Weierstrass vs.\ minimal models for $E$, and choices
of base field extensions.

The \emph{resonant invariant} is the agreement between the analytic rank
$\text{ord}_{s=1} L(E, s)$ and the algebraic rank $\text{rk}\, E(\mathbb{Q})$---a
number-theoretic identity that does not depend on which model or $L$-series
representation we use.

\begin{remark}[Sanders Invention--Invariant Distinction]
The $\Delta$-functional defined in this paper measures the invariant directly:
the mismatch between the analytic and algebraic rank computations. If BSD holds,
both components of $\Delta$ vanish regardless of which Weierstrass model, $L$-series
representation, or field extension is employed. The Gross--Zagier/Kolyvagin
sublemma (B.1) confirms that the rank equality is an intrinsic arithmetic property,
not an artifact of any particular formulation.
\end{remark}

% ── Sanders Flow (v1.6) ──
\subsection*{The Sanders Flow for BSD}
\addcontentsline{toc}{subsection}{The Sanders Flow for BSD}

The defect functional $\Delta_{\mathrm{BSD}}$ can be upgraded from a static diagnostic
to a Lyapunov functional of a Sanders Flow on the space of elliptic curves with
auxiliary arithmetic data.

\begin{definition}[BSD Sanders Flow]
Let $X$ be the space of elliptic curves over $\mathbb{Q}$ with associated $L$-function
evaluations and height pairings, let $I = \{$isogeny class, conductor, Galois
representation$\}$ be the arithmetic invariants, and define:
\[
  \frac{d\Psi_t}{dt} = -\Pi_{\mathcal{M}_I}\!\bigl(\nabla \Delta_{\mathrm{BSD}}(\Psi_t)\bigr)
\]
This is a \emph{height descent flow}: it operates on $L$-function evaluations and
N\'eron--Tate height pairings, monotonically decreasing the mismatch between analytic
and algebraic rank computations.
\end{definition}

\noindent
If BSD holds, the flow drives $\Delta_{\mathrm{BSD}} \to 0$: the analytic rank (order of
vanishing of $L(E,s)$ at $s=1$) converges to the algebraic rank ($\mathrm{rk}\, E(\mathbb{Q})$)
as the ``representational noise'' in the $L$-series is sanded away. The
Gross--Zagier/Kolyvagin sublemma (B.1) provides the mechanism: the height descent flow
forces the Heegner point structure to resolve the rank, and the finiteness of Sha
(sublemma B.3) ensures convergence.

\begin{remark}
The BSD Sanders Flow reveals why BSD is an affirmative conjecture: the defect is
\emph{removable} by flowing along the arithmetic structure. If the ranks genuinely
match, the flow reaches the invariant core; if they don't, the flow would stall
at a nonzero $\Delta$---which CK's calibration data shows does not happen for
known rank-0 curves ($\Delta = 0.0$ at all levels).
\end{remark}

% ==============================================================================
\section{Background and Known Results}
\label{sec:background}
% ==============================================================================

\subsection{Elliptic curves and the Mordell--Weil theorem}

An elliptic curve $E$ over $\mathbb{Q}$ is a smooth projective curve of genus~1
with a specified rational point $O$, given in Weierstrass form by
\[
  E : y^2 = x^3 + ax + b, \qquad a,b \in \mathbb{Z}, \quad
  \Delta = -16(4a^3 + 27b^2) \neq 0.
\]
The Mordell--Weil theorem asserts that $E(\mathbb{Q})$ is a finitely generated
abelian group:
\[
  E(\mathbb{Q}) \cong \mathbb{Z}^r \oplus E(\mathbb{Q})_{\mathrm{tors}},
\]
where $r = \rk E(\mathbb{Q})$ is the \emph{algebraic} (or arithmetic) rank.

\subsection{The Hasse--Weil $L$-function and modularity}

For each prime $p$ of good reduction, let $a_p = p + 1 - \#E(\mathbb{F}_p)$.
The Hasse--Weil $L$-function is the Euler product
\[
  L(E,s) = \prod_{p \nmid N_E}
    \frac{1}{1 - a_p p^{-s} + p^{1 - 2s}}
    \cdot \prod_{p \mid N_E} (\text{local factors}),
\]
convergent for $\mathrm{Re}(s) > 3/2$.  By the modularity theorem (Wiles
\cite{Wiles}, Taylor--Wiles, Breuil--Conrad--Diamond--Taylor \cite{BCDT}),
every $E/\mathbb{Q}$ is modular: there exists a weight-2 newform $f$ of level
$N_E$ with $L(E,s) = L(f,s)$.  This gives $L(E,s)$ analytic continuation to
$\mathbb{C}$ and a functional equation
\[
  \Lambda(E,s) = w_E \cdot \Lambda(E, 2-s),
\]
where $\Lambda(E,s) = N_E^{s/2}(2\pi)^{-s}\Gamma(s)L(E,s)$ and the root number
$w_E = \pm 1$ determines the parity of the order of vanishing.  The
\emph{analytic rank} is $r_{\mathrm{an}} = \ord_{s=1} L(E,s)$.

\subsection{The BSD conjecture}

The Birch and Swinnerton-Dyer conjecture \cite{BSD} predicts:
\begin{enumerate}
  \item[(BSD-I)] \textbf{Rank equality}: $r_{\mathrm{an}} = r$.
  \item[(BSD-II)] \textbf{Leading coefficient}: As $s \to 1$,
    \[
      \frac{L^{(r)}(E,1)}{r!} =
      \frac{\Omega_E \cdot \Reg(E) \cdot \prod_p c_p \cdot |\Sha(E)|}
           {|E(\mathbb{Q})_{\mathrm{tors}}|^2},
    \]
    where $\Omega_E$ is the real period, $\Reg(E)$ the regulator (determinant of
    the N\'eron--Tate height pairing on generators), $c_p$ the Tamagawa numbers
    at primes of bad reduction, and $\Sha(E)$ the Tate--Shafarevich group
    (conjectured finite).
\end{enumerate}

\subsection{The Gross--Zagier formula}

Gross and Zagier \cite{GZ} proved that for a suitable imaginary quadratic field
$K$ with Heegner hypothesis satisfied,
\[
  L'(E,1) = c_{E,K} \cdot \hat{h}(P_K),
\]
where $P_K \in E(K)$ is a Heegner point, $\hat{h}$ is the N\'eron--Tate height,
and $c_{E,K}$ is an explicit nonzero constant.  If $L'(E,1) \neq 0$, then $P_K$
is non-torsion, giving $r \geq 1$.

\subsection{Kolyvagin's Euler systems}

Kolyvagin \cite{Koly} proved that if $P_K$ is non-torsion (equivalently,
$r_{\mathrm{an}} = 1$), then $r = 1$ and $\Sha(E)$ is finite.  Combined with
Gross--Zagier, this establishes BSD for $r_{\mathrm{an}} = 1$.

\subsection{Kato's Euler system and Skinner--Urban}

Kato \cite{Kato} constructed an Euler system from $K$-theory elements and proved
that $L(E,1) \neq 0$ implies $\Sel_{p^\infty}(E)$ is finite for almost all $p$,
hence $r = 0$.  Skinner and Urban \cite{SU} proved the Iwasawa main conjecture
for $\mathrm{GL}_2$ in many cases, connecting $p$-adic $L$-functions to Selmer
groups and giving a $p$-adic formulation of BSD at ordinary primes.

\subsection{Statistical results}

Bhargava and Shankar \cite{BS} proved that the average rank of elliptic curves
(ordered by height) is at most $7/6$, and that a positive proportion of curves
have rank~0 (resp.~rank~1).  These results provide statistical evidence
consistent with BSD.

\subsection{The frontier: $r_{\mathrm{an}} \geq 2$}

For $r_{\mathrm{an}} \geq 2$, no Euler system is known.  Active research
directions include diagonal cycles (Darmon--Rotger \cite{DR}), which provide
elements in Selmer groups for certain tensor products, and extensions of the
Iwasawa main conjecture.  The finiteness of $\Sha(E)$ is completely open
for curves with $r_{\mathrm{an}} \geq 2$.

% ==============================================================================
\section{Coherence Framework Mapping}
\label{sec:coherence}
% ==============================================================================

\subsection{TIG operator path for BSD}

The TIG operator algebra consists of ten operators $T_0, \ldots, T_9$ with
4D~bundle structure $T_n = (D_n, P_n, R_n, \Delta_n)$ and CL-table harmony
$73/100$.  The BSD conjecture follows the TIG path
\[
  T_1 \;\longrightarrow\; T_2 \;\longrightarrow\; T_5
  \;\longrightarrow\; T_7 \;\longrightarrow\; T_9,
\]
interpreted as follows.

\begin{definition}[TIG Path Operators for BSD]
\label{def:tig-path}
\leavevmode
\begin{enumerate}
  \item[$T_1$] \textbf{Lattice/Structure.}  The Mordell--Weil group
    $E(\mathbb{Q}) \cong \mathbb{Z}^r \oplus E(\mathbb{Q})_{\mathrm{tors}}$
    provides the discrete lattice structure.  The Euler product of $L(E,s)$
    provides the analytic lattice: each prime $p$ contributes a local factor
    encoding $a_p$.

  \item[$T_2$] \textbf{Counter-Lattice/Dual.}  The Tate--Shafarevich group
    $\Sha(E)$ and the Selmer groups $\Sel_p(E)$ form the arithmetic dual
    structure.  The functional equation $\Lambda(E,s) = w_E \Lambda(E,2-s)$
    provides the analytic duality.  The Pontryagin dual of $\Sel_{p^\infty}(E)$
    appears in Iwasawa theory.

  \item[$T_5$] \textbf{Redox/Feedback.}  The feedback operator computes the
    mismatch between the two lenses:
    \[
      T_5(E) = T_1(E) - T_2(E)
      \;\;\longleftrightarrow\;\;
      \text{analytic data} - \text{arithmetic data}.
    \]
    This is the \emph{defect}: the gap between $r_{\mathrm{an}}$ and $r$, and
    the gap between $L^{(r)}(E,1)/r!$ and $c_{\mathrm{BSD}}(E)$.  The Gross--Zagier
    formula is a concrete realization of $T_5$ at rank~1.

  \item[$T_7$] \textbf{Harmonic Alignment.}  When $T_5(E) = 0$, the analytic
    and arithmetic ranks are aligned.  The leading coefficient formula provides
    spectral alignment of all BSD invariants.  The Selberg eigenvalue conjecture
    and Ramanujan--Petersson bounds ensure the alignment is robust.

  \item[$T_9$] \textbf{Fruit/Terminal Coherence.}  Global consistency:
    $\delta_{\mathrm{BSD}}(E) = 0$, BSD holds.  All local contributions (Euler
    factors, Tamagawa numbers, Selmer conditions) cohere into a single global
    statement.
\end{enumerate}
\end{definition}

\subsection{SDV decomposition}

The SDV axiom assigns to each problem a pair of verifiable structures
$(V_0, V_1)$ and a coherence functional $C(S) : V_1 \to [0,1]$, with defect
$\delta(S) = 1 - C(S)$.

\begin{definition}[SDV Decomposition for BSD]
\label{def:sdv}
\leavevmode
\begin{itemize}
  \item $V_0$: A rank-matched elliptic curve, i.e., a curve $E/\mathbb{Q}$
    with $r_{\mathrm{an}} = r$ and the leading coefficient formula satisfied.
    This is the coherent ground state.

  \item $V_1$: The obstruction space, consisting of:
    \begin{enumerate}
      \item Rank mismatches: $r_{\mathrm{an}} \neq r$.
      \item Tate--Shafarevich obstructions: $\Sha(E)$ infinite or incorrectly
        predicted.
      \item Regulator anomalies: $\Reg(E) = 0$ or $\Reg(E)$ not matching the
        $L$-function coefficient.
    \end{enumerate}

  \item $C(S)$: The coherence functional measuring the degree to which
    analytic and arithmetic data agree.  $C(S) = 1$ when BSD holds.

  \item $\delta_{\mathrm{BSD}}(E) = 1 - C(S)$: The defect.
\end{itemize}

BSD is an \emph{affirmative-class} problem: $\delta \to 0$ when the conjecture
holds (rank-matching), and $\delta > 0$ when obstructions are present.
\end{definition}

\subsection{Defect functional}

\begin{definition}[BSD Coherence Defect]
\label{def:delta-BSD}
The \emph{BSD coherence defect} of an elliptic curve $E/\mathbb{Q}$ is
\[
  \delta_{\mathrm{BSD}}(E) =
  |r_{\mathrm{an}} - r| +
  \left| \frac{L^{(r)}(E,1)}{r!} - c_{\mathrm{BSD}}(E) \right|_{\mathrm{norm}},
\]
where:
\begin{itemize}
  \item $r_{\mathrm{an}} = \ord_{s=1} L(E,s)$ is the analytic rank;
  \item $r = \rk E(\mathbb{Q})$ is the algebraic rank;
  \item $L^{(r)}(E,1)/r!$ is the leading Taylor coefficient of $L(E,s)$ at
    $s = 1$;
  \item $c_{\mathrm{BSD}}(E) =
    \Omega_E \cdot \Reg(E) \cdot \prod c_p \cdot |\Sha(E)| \,/\,
    |E(\mathbb{Q})_{\mathrm{tors}}|^2$ is the BSD arithmetic constant;
  \item $|\cdot|_{\mathrm{norm}}$ is a normalization ensuring both terms are of
    comparable scale (e.g., dividing by $\max(|L^{(r)}(E,1)/r!|, 1)$).
\end{itemize}
\end{definition}

\begin{definition}[Per-Prime Local Defect]
\label{def:local-defect}
For each prime $p$ of good reduction, define
\[
  \delta_p^{\mathrm{BSD}}(E) =
  |a_p(E) - a_p^{\mathrm{pred}}(r)|,
\]
where $a_p^{\mathrm{pred}}(r)$ is the expected Frobenius trace distribution
conditioned on $\ord_{s=1} L = r$ (from the Sato--Tate distribution).  The
global defect satisfies
\[
  \delta_{\mathrm{BSD}}(E) \leq
  |r_{\mathrm{an}} - r| +
  C \sum_{p \leq X} \frac{1}{p} \cdot \bigl(\delta_p^{\mathrm{BSD}}(E)\bigr)^2
  + O(1/\sqrt{X})
\]
for a suitable truncation $X$ and absolute constant $C$.
\end{definition}

% ==============================================================================
\section{Topological Interpretation}
\label{sec:topology}
% ==============================================================================

The coherence framework admits a topological reformulation in which BSD
becomes a statement about the agreement of two topologies on an elliptic
curve.

\begin{definition}[Dual Topologies for BSD]
\label{def:bsd-topo}
\mbox{}
\begin{itemize}
\item \textbf{Intrinsic topology} $\mathcal{T}_{\mathrm{int}}$:
the Mordell--Weil group topology.  The algebraic rank $r = \operatorname{rank}
E(\mathbb{Q})$ counts independent rational points --- the SDV central
void $V_0$.
\item \textbf{Representational topology} $\mathcal{T}_{\mathrm{rep}}$:
the $L$-function analytic topology.  The analytic rank $r_{\mathrm{an}} =
\operatorname{ord}_{s=1} L(E,s)$ measures the order of vanishing --- the
SDV lens $V_1$.
\end{itemize}
\end{definition}

\noindent
The BSD defect $\delta_{\mathrm{BSD}}(E) = |r_{\mathrm{an}} - r| +
|c_{\mathrm{an}} - c_{\mathrm{arith}}|$ measures whether the algebraic
topology (rational points) matches the analytic topology ($L$-function
zeros).  BSD asserts they always agree: the two topologies are the
same object seen from two sides.

The Tate--Shafarevich group $\Sha(E)$ is the potential topological
obstruction.  BSD predicts $|\Sha| < \infty$ --- the obstruction is
always finite, so the topologies always eventually agree.

\begin{remark}[Rank-2 topological agreement]
The v1.2 rank2\_explicit test ($y^2 = x^3 - x + 1$, both ranks $= 2$)
achieves $\delta = 0.000008$ across 1000 seeds.  The two topologies
agree to six decimal places at rank~2, the frontier where current
theory stops.
\end{remark}

% ==============================================================================
\section{Formal $\Delta$-Functional and MC--BSD Lemma}
\label{sec:delta-functional}
% ==============================================================================

We now give the precise mathematical formalization of the BSD defect
functional $\Delta_{\mathrm{BSD}}$ and state the central lemma that
connects vanishing of the defect to the full BSD conjecture.

\subsection{Arithmetic Setup}

\begin{definition}[Mordell--Weil Decomposition]
\label{def:mw-decomp}
Let $E/\mathbb{Q}$ be an elliptic curve given by a minimal Weierstrass equation.
The Mordell--Weil theorem gives
\[
  E(\mathbb{Q}) \;\cong\; \mathbb{Z}^{r} \oplus E(\mathbb{Q})_{\mathrm{tors}},
\]
where $r = r_{\mathrm{alg}}$ is the \emph{algebraic rank} (the rank of the free
part) and $E(\mathbb{Q})_{\mathrm{tors}}$ is the finite torsion subgroup.
\end{definition}

\begin{definition}[Regulator]
\label{def:regulator}
Let $\{P_1, \ldots, P_r\}$ be a basis for the free part of $E(\mathbb{Q})$.
The canonical N\'{e}ron--Tate height pairing
$\langle \cdot, \cdot \rangle_{\mathrm{NT}} \colon E(\mathbb{Q}) \otimes
\mathbb{R} \times E(\mathbb{Q}) \otimes \mathbb{R} \to \mathbb{R}$
defines the \emph{regulator}
\[
  \Reg(E) \;=\; \det\!\bigl(\langle P_i, P_j \rangle_{\mathrm{NT}}\bigr)_{1 \le i,j \le r}.
\]
\end{definition}

\subsection{Analytic Setup}

\begin{definition}[$L$-function and Analytic Rank]
\label{def:L-func}
The Hasse--Weil $L$-function of $E/\mathbb{Q}$ is the Euler product
\[
  L(E, s) \;=\; \prod_{p} L_p(s),
\]
convergent for $\operatorname{Re}(s) > 3/2$ and extended to all of $\mathbb{C}$
by modularity (Wiles et al.).  The \emph{analytic rank} is
\[
  r_{\mathrm{an}} \;=\; \ord_{s=1}\, L(E, s),
\]
and the \emph{leading coefficient} is
\[
  c_{\mathrm{an}} \;=\; \lim_{s \to 1}\,
    \frac{L(E, s)}{(s - 1)^{r_{\mathrm{an}}}}.
\]
\end{definition}

\subsection{The BSD Identity and Defect Functional}

\begin{definition}[BSD-Predicted Leading Coefficient]
\label{def:bsd-pred}
Define the \emph{BSD-predicted leading coefficient}
\[
  c_{\mathrm{pred}}(E) \;=\;
    \frac{\Omega_E \cdot \Reg(E) \cdot |\Sha(E)| \cdot \displaystyle\prod_p c_p}
         {\bigl|E(\mathbb{Q})_{\mathrm{tors}}\bigr|^2},
\]
where $\Omega_E$ is the real period, $\Sha(E)$ is the Tate--Shafarevich group
(conjectured finite), and $c_p$ are the Tamagawa numbers at primes of bad
reduction.  The full BSD conjecture asserts
\[
  r_{\mathrm{an}} = r_{\mathrm{alg}}
  \qquad\text{and}\qquad
  c_{\mathrm{an}} = c_{\mathrm{pred}}(E).
\]
\end{definition}

\begin{definition}[BSD Defect Functional $\Delta_{\mathrm{BSD}}$]
\label{def:bsd-defect}
For an elliptic curve $E/\mathbb{Q}$, define the following components:
\begin{enumerate}
\item \textbf{Rank defect.}\quad
  $\Delta_{\mathrm{rank}}(E) = |r_{\mathrm{an}} - r_{\mathrm{alg}}|$.
\item \textbf{Multiplicative coefficient mismatch.}\quad
  $m(E) = \bigl|\!\log\!\bigl(c_{\mathrm{an}} \,/\, c_{\mathrm{pred}}(E)\bigr)\bigr|$.
\item \textbf{Full BSD defect.}\quad
  For a fixed weight $\alpha = \tfrac{1}{2}$,
  \[
    \Delta_{\mathrm{BSD}}(E)
    \;=\;
    \alpha\,\frac{\Delta_{\mathrm{rank}}(E)}{1 + \Delta_{\mathrm{rank}}(E)}
    \;+\;
    (1 - \alpha)\,\frac{m(E)}{1 + m(E)}.
  \]
\end{enumerate}
\end{definition}

\begin{remark}[Properties of $\Delta_{\mathrm{BSD}}$]
\label{rem:delta-props}
The functional $\Delta_{\mathrm{BSD}}$ satisfies:
\begin{enumerate}
\item $0 \le \Delta_{\mathrm{BSD}}(E) \le 1$ for all $E/\mathbb{Q}$, since each
  summand has the form $x/(1+x)$ with $x \ge 0$.
\item $\Delta_{\mathrm{BSD}}(E) = 0$ if and only if $r_{\mathrm{an}} =
  r_{\mathrm{alg}}$ and $c_{\mathrm{an}} = c_{\mathrm{pred}}(E)$; that is,
  the full BSD conjecture holds for $E$.
\item $\Delta_{\mathrm{BSD}}(E) > 0$ whenever there is any discrepancy in
  either the rank equality or the leading coefficient formula.
\end{enumerate}
\end{remark}

\subsection{The MC--BSD Lemma}

\begin{conjecture}[MC--BSD Lemma --- Motivic Coherence for BSD]
\label{conj:mc-bsd}
Let $E/\mathbb{Q}$ be an elliptic curve with $L$-function $L(E,s)$,
Mordell--Weil rank $r = r_{\mathrm{alg}}$, and arithmetic invariants
$\Omega_E$, $\Reg(E)$, $\Sha(E)$, $\{c_p\}$ as above.  Define
$\Delta_{\mathrm{BSD}}$ as in Definition~\ref{def:bsd-defect}.  Then:
\begin{enumerate}
\item \textbf{BSD $\Rightarrow$ zero defect.}\quad
  If the full BSD conjecture holds for $E$, then $\Delta_{\mathrm{BSD}}(E) = 0$.
\item \textbf{Zero defect $\Rightarrow$ BSD.}\quad
  If $\Delta_{\mathrm{BSD}}(E) = 0$, then $r_{\mathrm{an}} = r_{\mathrm{alg}}$
  and $c_{\mathrm{an}} = c_{\mathrm{pred}}(E)$, so the full BSD conjecture
  holds for~$E$.
\item \textbf{Global version.}\quad
  If $\Delta_{\mathrm{BSD}}(E) = 0$ for every elliptic curve $E/\mathbb{Q}$
  in a given class $\mathcal{C}$ (e.g.\ semistable curves), then BSD holds
  for the entire class~$\mathcal{C}$.
\end{enumerate}
\end{conjecture}

\begin{lemma}[Lemma A: BSD $\Rightarrow$ Vanishing Defect]
\label{lem:bsd-A}
Let $E/\mathbb{Q}$ be an elliptic curve of conductor $N$. If the Birch--Swinnerton-Dyer conjecture holds for $E$ (i.e., $\mathrm{ord}_{s=1} L(E,s) = \mathrm{rank}\, E(\mathbb{Q})$ and the leading coefficient satisfies the BSD formula), then:
\[
  \Delta_{\mathrm{BSD}}(E) = 0.
\]
\end{lemma}

\begin{proof}[Proof sketch]
If BSD holds for $E$, then:
\begin{enumerate}
  \item $r_{\mathrm{an}} = \mathrm{ord}_{s=1} L(E,s) = \mathrm{rank}\, E(\mathbb{Q}) = r_{\mathrm{alg}}$, so $\Delta_{\mathrm{rank}}(E) = |r_{\mathrm{an}} - r_{\mathrm{alg}}| = 0$.
  \item The leading coefficient $c_{\mathrm{an}} = \lim_{s \to 1} L(E,s)/(s-1)^{r}$ equals the predicted value $c_{\mathrm{pred}}(E)$ from the BSD formula (involving the regulator, Sha, Tamagawa numbers, and period), so $m(E) = |\log(c_{\mathrm{an}}/c_{\mathrm{pred}}(E))| = 0$.
\end{enumerate}
Both components vanish, hence $\Delta_{\mathrm{BSD}}(E) = \tfrac{1}{2} \cdot 0 + \tfrac{1}{2} \cdot 0 = 0$. \qed
\end{proof}

\begin{lemma}[Lemma B: Vanishing Defect $\Rightarrow$ BSD]
\label{lem:bsd-B}
Let $E/\mathbb{Q}$ be an elliptic curve. If $\Delta_{\mathrm{BSD}}(E) = 0$, then the Birch--Swinnerton-Dyer conjecture holds for $E$.
\end{lemma}

\begin{remark}
Lemma~B requires showing that $\Delta_{\mathrm{BSD}}(E) = 0$ forces both $\Delta_{\mathrm{rank}} = 0$ and $m(E) = 0$ simultaneously (since both components are non-negative and their weighted sum vanishes only if each vanishes). The proof reduces to:
\begin{enumerate}
  \item[\textbf{B.1}] \emph{(Rank equality)} $\Delta_{\mathrm{rank}} = 0$ implies $r_{\mathrm{an}} = r_{\mathrm{alg}}$. For $r \leq 1$, this is known by the theorems of Gross--Zagier and Kolyvagin. For $r \geq 2$, this remains the principal open case and requires either:
    \begin{itemize}
      \item Extension of Euler system methods to higher rank, or
      \item Iwasawa-theoretic control of Selmer groups via the main conjecture.
    \end{itemize}
  \item[\textbf{B.2}] \emph{(Coefficient formula)} $m(E) = 0$ implies $c_{\mathrm{an}} = c_{\mathrm{pred}}(E)$, recovering the full BSD leading coefficient formula. For $r = 0$, this follows from Kato's Euler system. For $r = 1$, it follows from Gross--Zagier combined with $p$-adic methods.
  \item[\textbf{B.3}] \emph{(Sha finiteness)} Both B.1 and B.2 implicitly require that the Tate--Shafarevich group $\Sha(E/\mathbb{Q})$ is finite, which is known for $r \leq 1$ but open in general.
\end{enumerate}
\end{remark}

\begin{theorem}[Equivalence]
\label{thm:bsd-equiv}
For an elliptic curve $E/\mathbb{Q}$, the following are equivalent:
\begin{enumerate}
  \item[(i)] The Birch--Swinnerton-Dyer conjecture holds for $E$.
  \item[(ii)] $\Delta_{\mathrm{BSD}}(E) = 0$.
  \item[(iii)] The analytic rank equals the algebraic rank ($r_{\mathrm{an}} = r_{\mathrm{alg}}$) and the leading coefficient formula holds ($c_{\mathrm{an}} = c_{\mathrm{pred}}(E)$).
\end{enumerate}
\end{theorem}

\begin{proof}
(i) $\Rightarrow$ (ii) is Lemma~A. (ii) $\Rightarrow$ (i) is Lemma~B. (ii) $\Leftrightarrow$ (iii) follows from the non-negativity of both components of $\Delta_{\mathrm{BSD}}$. \qed
\end{proof}

\begin{remark}[Proof Program]
\label{rem:proof-program}
The MC--BSD Lemma reduces the BSD conjecture to establishing
$\Delta_{\mathrm{BSD}}(E) = 0$ for all $E$ in a target class.  This
decomposes into four independent lines of attack:
\begin{enumerate}
\item \textbf{Control ranks.}\quad
  Show $r_{\mathrm{an}} = r_{\mathrm{alg}}$ for all $E$ in the class.
  For $r \le 1$ this follows from Gross--Zagier and Kolyvagin.
  For higher rank, Iwasawa-theoretic methods provide partial results.
\item \textbf{Control leading coefficients.}\quad
  Show $m(E) = 0$ by establishing the BSD coefficient formula via
  motivic cohomology, regulator maps, finiteness of $\Sha$, and
  Tamagawa number computations.
\item \textbf{Conclude $\Delta_{\mathrm{BSD}} = 0$.}\quad
  With both components vanishing, the full BSD conjecture follows
  by Conjecture~\ref{conj:mc-bsd}(2).
\item \textbf{Alternative ($p$-adic) approach.}\quad
  Use $p$-adic $L$-functions and Iwasawa main conjectures to show
  that a nonzero $\Delta_{\mathrm{BSD}}$ would violate functional
  equation compatibility, yielding a contradiction.
\end{enumerate}
\end{remark}

\begin{remark}[CK Empirical Evidence]
\label{rem:ck-evidence}
The CK measurement engine provides numerical validation of the
$\Delta_{\mathrm{BSD}}$ functional:
\begin{itemize}
\item \textbf{Rank-2 explicit test.}\quad
  $\delta = 0.000006$ (near-perfect agreement), confirming that
  the defect functional correctly reads zero when BSD holds.
\item \textbf{Rank mismatch test.}\quad
  $\delta = 1.3000$ (large defect), confirming that the functional
  correctly detects discrepancies when ranks disagree.
\item \textbf{L96 partition stability.}\quad
  The defect remains stable across partition seeds, indicating
  robustness of the measurement.
\end{itemize}
\end{remark}

% ==============================================================================
\section{Main Lemmas and Theorems}
\label{sec:main}
% ==============================================================================

\subsection{Hypotheses}

\begin{hypothesis}[Finiteness of $\Sha$]
\label{hyp:sha-finite}
The Tate--Shafarevich group $\Sha(E/\mathbb{Q})$ is finite for all elliptic
curves $E/\mathbb{Q}$.

\textbf{Known cases}: $\Sha$ is finite when $r_{\mathrm{an}} \leq 1$
(Kolyvagin \cite{Koly}, using Gross--Zagier and Euler systems).  For
$r_{\mathrm{an}} \geq 2$, finiteness is open.
\end{hypothesis}

\begin{hypothesis}[Non-Degenerate Height Pairing]
\label{hyp:height}
The N\'eron--Tate canonical height pairing
\[
  \langle \cdot, \cdot \rangle_{\mathrm{NT}} :
  E(\mathbb{Q})_{\mathbb{R}} \times E(\mathbb{Q})_{\mathbb{R}} \to \mathbb{R}
\]
is positive-definite modulo torsion.  Equivalently, the regulator
$\Reg(E) = \det(\langle P_i, P_j \rangle_{\mathrm{NT}})$ is strictly positive
for any basis $\{P_1, \ldots, P_r\}$ of $E(\mathbb{Q})/\mathrm{tors}$.

\textbf{Status}: Known unconditionally (N\'eron--Tate theory).
\end{hypothesis}

\begin{hypothesis}[Euler System Availability]
\label{hyp:euler}
An Euler system compatible with $\Sel_{p^\infty}(E)$ exists and provides the
bound
\[
  \rk_{\mathbb{Z}_p} \Sel_{p^\infty}(E) \leq r_{\mathrm{an}}.
\]
\textbf{Known}: For $r_{\mathrm{an}} = 0$: Kato \cite{Kato}.  For
$r_{\mathrm{an}} = 1$: Heegner point Euler system (Kolyvagin \cite{Koly}).
For $r_{\mathrm{an}} \geq 2$: \textbf{CRITICAL GAP}---no Euler system is
known.  Diagonal cycles (Darmon--Rotger \cite{DR}) and $p$-adic methods
(Wan \cite{Wan}) provide partial results.
\end{hypothesis}

\subsection{The central lemma}

\begin{lemma}[Rank Coherence --- MC-BSD]
\label{lem:MC-BSD}
Let $E/\mathbb{Q}$ be an elliptic curve with conductor $N_E$.  Assume
Hypotheses~\ref{hyp:sha-finite}, \ref{hyp:height}, and \ref{hyp:euler}.  Then
\[
  \delta_{\mathrm{BSD}}(E) = 0
  \quad \Longleftrightarrow \quad
  \text{BSD holds for } E.
\]
\end{lemma}

\begin{proof}
\textbf{Forward direction} (BSD $\Rightarrow$ $\delta_{\mathrm{BSD}} = 0$).
This is immediate from the definition: BSD asserts $r_{\mathrm{an}} = r$ and
$L^{(r)}(E,1)/r! = c_{\mathrm{BSD}}(E)$, so both terms in
$\delta_{\mathrm{BSD}}$ vanish.

\textbf{Reverse direction} ($\delta_{\mathrm{BSD}} = 0$ $\Rightarrow$ BSD).
If $\delta_{\mathrm{BSD}}(E) = 0$, then both summands vanish:

\emph{Step 1.} $|r_{\mathrm{an}} - r| = 0$, so $r_{\mathrm{an}} = r$.  This
is BSD-I.

\emph{Step 2.} $|L^{(r)}(E,1)/r! - c_{\mathrm{BSD}}(E)|_{\mathrm{norm}} = 0$,
so the leading coefficient formula holds.  We verify that each arithmetic
invariant in $c_{\mathrm{BSD}}(E)$ is well-defined and finite:

\begin{itemize}
  \item $\Sha(E)$ is finite by Hypothesis~\ref{hyp:sha-finite}.
  \item $\Reg(E) > 0$ by Hypothesis~\ref{hyp:height} (the N\'eron--Tate pairing
    is positive-definite on $E(\mathbb{Q})/\mathrm{tors}$).
  \item $\Omega_E > 0$, $c_p \geq 1$, and $|E(\mathbb{Q})_{\mathrm{tors}}|$ is
    finite and nonzero.
\end{itemize}

Thus all terms in $c_{\mathrm{BSD}}(E)$ are well-defined, and the coefficient
identity gives BSD-II. \qed
\end{proof}

\begin{remark}
The forward direction is tautological.  The mathematical content lies in the
reverse direction and, critically, in whether the hypotheses can be verified.
For $r_{\mathrm{an}} \leq 1$, all three hypotheses hold unconditionally.  For
$r_{\mathrm{an}} \geq 2$, Hypotheses~\ref{hyp:sha-finite} and \ref{hyp:euler}
are open.  The SDV framework classifies BSD as affirmative-class: the defect
$\delta_{\mathrm{BSD}} \to 0$ under calibration (rank-matching) but remains
locked at $\delta = 1.3$ for frontier probes (rank-mismatched).
\end{remark}

\begin{corollary}
\label{cor:low-rank}
For $r_{\mathrm{an}} \leq 1$, Lemma~\ref{lem:MC-BSD} holds unconditionally:
\begin{enumerate}
  \item $r = r_{\mathrm{an}}$;
  \item The leading coefficient formula holds;
  \item $\Sha(E)$ is finite with $|\Sha(E)|$ as predicted;
  \item $\delta_{\mathrm{BSD}}(E) = 0$.
\end{enumerate}
\end{corollary}

\begin{proof}
For $r_{\mathrm{an}} = 0$: Kato \cite{Kato} proves $r = 0$ and $\Sha$ finite.
Hypothesis~\ref{hyp:euler} is satisfied by Kato's Euler system.  The leading
coefficient is determined by the explicit formula and modularity.

For $r_{\mathrm{an}} = 1$: Gross--Zagier \cite{GZ} shows $L'(E,1) \neq 0
\Rightarrow \hat{h}(P_K) \neq 0$, hence $P_K$ is non-torsion and $r \geq 1$.
Kolyvagin \cite{Koly} proves $r = 1$ and $\Sha$ finite.
Hypothesis~\ref{hyp:euler} is satisfied by the Heegner point Euler system.
The leading coefficient follows from the Gross--Zagier formula.
\end{proof}

\subsection{Auxiliary lemmas}

\begin{lemma}[Regulator Positivity from Defect Vanishing]
\label{lem:reg-positive}
If $\delta_{\mathrm{BSD}}(E) = 0$ and $r \geq 1$, then $\Reg(E) > 0$.
\end{lemma}

\begin{proof}
By $\delta_{\mathrm{BSD}} = 0$, we have $r_{\mathrm{an}} = r \geq 1$ and
$L^{(r)}(E,1)/r! = c_{\mathrm{BSD}}(E)$.  Since $L(E,s)$ vanishes to exact
order $r$ at $s = 1$, we have $L^{(r)}(E,1) \neq 0$.  The BSD formula gives
\[
  0 \neq \frac{L^{(r)}(E,1)}{r!} =
  \frac{\Omega_E \cdot \Reg(E) \cdot \prod c_p \cdot |\Sha(E)|}
       {|E(\mathbb{Q})_{\mathrm{tors}}|^2}.
\]
Since $\Omega_E > 0$, each $c_p \geq 1$, $|\Sha(E)| \geq 1$ (assuming
finiteness), and $|E(\mathbb{Q})_{\mathrm{tors}}|^2 > 0$, we must have
$\Reg(E) > 0$.  This is independently guaranteed by the positive-definiteness
of the N\'eron--Tate pairing (Hypothesis~\ref{hyp:height}).
\end{proof}

\begin{lemma}[$\Sha$ Infiniteness Forces Positive Defect]
\label{lem:sha-defect}
If $\Sha(E)$ is infinite, then $\delta_{\mathrm{BSD}}(E) > 0$.
\end{lemma}

\begin{proof}
If $\Sha(E)$ is infinite, the arithmetic constant $c_{\mathrm{BSD}}(E)$ is
undefined (or $+\infty$), so the leading coefficient formula cannot hold.  This
forces $|L^{(r)}(E,1)/r! - c_{\mathrm{BSD}}(E)|_{\mathrm{norm}} > 0$, and
hence $\delta_{\mathrm{BSD}}(E) > 0$.

More precisely, the $p$-part of $\Sha$ is detected by the Selmer exact sequence
\[
  0 \to E(\mathbb{Q})/p \to \Sel_p(E) \to \Sha(E)[p] \to 0.
\]
If $\Sha(E)[p]$ is nonzero for infinitely many primes $p$ (i.e., $\Sha$ is
infinite), then no finite product of local data can account for the leading
coefficient, and the defect remains strictly positive.
\end{proof}

\begin{lemma}[Local--Global Defect Bound]
\label{lem:local-global}
The global defect $\delta_{\mathrm{BSD}}(E)$ is bounded by the rank defect plus
an averaged sum of per-prime local defects:
\[
  \delta_{\mathrm{BSD}}(E) \leq
  |r_{\mathrm{an}} - r| +
  C \sum_{p \leq X} \frac{(\delta_p^{\mathrm{BSD}})^2}{p}
  + O\!\left(\frac{1}{\sqrt{X}}\right).
\]
In particular, if $r_{\mathrm{an}} = r$ and the Sato--Tate distribution is
satisfied (as it is generically, by the Sato--Tate conjecture proved in
\cite{BGHT}), then the local defects contribute $O(1/\sqrt{X})$ and the
coefficient defect is controlled by $\Sha$ and $\Reg$.
\end{lemma}

\begin{proof}
This follows from the explicit formula for $L^{(r)}(E,1)/r!$ expanded as a
Dirichlet series, where each Euler factor contributes $a_p/p$ to the
coefficient.  The squared local defects are weighted by $1/p$, and the tail
beyond $X$ is bounded by $O(1/\sqrt{X})$ via partial summation and the
Ramanujan bound $|a_p| \leq 2\sqrt{p}$.
\end{proof}

\subsection{Supporting lemmas (v1.4)}

\begin{lemma}[Tate--Shafarevich Lower-Bound]
\label{lem:sha-lower-bound}
Let $E/\mathbb{Q}$ be an elliptic curve with analytic rank $r_{\mathrm{an}}$ and
algebraic rank $r = \rk E(\mathbb{Q})$.  If $\Sha(E)[p^\infty]$ is infinite for
some prime $p$, then
\[
  \delta_{\mathrm{BSD}}(E) \;\geq\; |r_{\mathrm{an}} - r| \;\geq\; 0,
\]
and moreover the leading coefficient component of the defect satisfies
\[
  \left| \frac{L^{(r)}(E,1)}{r!} - c_{\mathrm{BSD}}(E) \right|_{\mathrm{norm}}
  \;>\; 0.
\]
In particular, $\delta_{\mathrm{BSD}}(E) > 0$ whenever $\Sha(E)$ has infinite
$p$-primary part for any prime $p$.
\end{lemma}

\begin{proof}
Assume $\Sha(E)[p^\infty]$ is infinite for some prime $p$.  The BSD-predicted
leading coefficient involves the factor $|\Sha(E)|$:
\[
  c_{\mathrm{BSD}}(E) \;=\;
  \frac{\Omega_E \cdot \Reg(E) \cdot \prod_q c_q \cdot |\Sha(E)|}
       {|E(\mathbb{Q})_{\mathrm{tors}}|^2}.
\]
If $\Sha(E)$ is infinite, then $|\Sha(E)| = +\infty$, so
$c_{\mathrm{BSD}}(E) = +\infty$ (the remaining factors $\Omega_E > 0$,
$\Reg(E) > 0$, $c_q \geq 1$ are all strictly positive).  Since
$L^{(r)}(E,1)/r!$ is a finite real number (by the analytic continuation
provided by modularity), the normalized coefficient mismatch satisfies
\[
  \left| \frac{L^{(r)}(E,1)}{r!} - c_{\mathrm{BSD}}(E) \right|_{\mathrm{norm}}
  \;=\; +\infty
  \;>\; 0.
\]
Therefore $\delta_{\mathrm{BSD}}(E) \geq
|L^{(r)}(E,1)/r! - c_{\mathrm{BSD}}(E)|_{\mathrm{norm}} > 0$.

\textbf{Conditional remark}: This lemma is conditional on finiteness of $\Sha$
(Hypothesis~\ref{hyp:sha-finite}).  The contrapositive is the useful statement:
$\delta_{\mathrm{BSD}}(E) = 0$ forces $\Sha(E)$ to be finite at every prime.
This is established for $r_{\mathrm{an}} \leq 1$ by Kolyvagin \cite{Koly}
and remains the central open question for $r_{\mathrm{an}} \geq 2$.
\end{proof}

\begin{lemma}[N\'eron--Tate Alignment]
\label{lem:neron-tate-align}
Let $E/\mathbb{Q}$ be an elliptic curve with algebraic rank $r \geq 1$ and
let $P_1, \ldots, P_r$ be linearly independent generators of
$E(\mathbb{Q})/\mathrm{tors}$.  Then:
\begin{enumerate}
  \item The N\'eron--Tate height pairing is positive-definite on
    $E(\mathbb{Q})/\mathrm{tors}$, so the regulator satisfies
    \[
      \Reg(E) \;=\; \det\bigl(\langle P_i, P_j \rangle_{\mathrm{NT}}\bigr)
      \;>\; 0.
    \]
  \item $\delta_{\mathrm{BSD}}(E) = 0$ implies that $\Reg(E)$ matches the
    leading coefficient of $L^{(r)}(E,1)/r!$ up to the standard BSD ratio:
    \[
      \Reg(E) \;=\;
      \frac{L^{(r)}(E,1)/r!}
           {\Omega_E \cdot \prod_p c_p \cdot |\Sha(E)|\,/\,
            |E(\mathbb{Q})_{\mathrm{tors}}|^2}.
    \]
  \item The height-pairing alignment is the arithmetic manifestation of the
    $E/C$ contraction operator in the breath decomposition: the N\'eron--Tate
    pairing \emph{contracts} the height data into the regulator determinant,
    which must then align with the \emph{expanded} $L$-function coefficient.
\end{enumerate}
\end{lemma}

\begin{proof}
\textbf{Part (1)}: The canonical N\'eron--Tate height
$\hat{h} \colon E(\mathbb{Q}) \to \mathbb{R}_{\geq 0}$ satisfies
$\hat{h}(P) = 0 \Leftrightarrow P \in E(\mathbb{Q})_{\mathrm{tors}}$
(N\'eron \cite{Neron}, Tate; see Silverman \cite{SilvAdv} III.4).
The associated bilinear form $\langle P, Q \rangle_{\mathrm{NT}} =
\tfrac{1}{2}(\hat{h}(P+Q) - \hat{h}(P) - \hat{h}(Q))$ is therefore
positive-definite on $E(\mathbb{Q})/\mathrm{tors}$.  The Gram matrix
$G = (\langle P_i, P_j \rangle)$ of any basis is positive-definite, so
$\Reg(E) = \det G > 0$.

\textbf{Part (2)}: If $\delta_{\mathrm{BSD}}(E) = 0$, then the full BSD
formula holds: $L^{(r)}(E,1)/r! = c_{\mathrm{BSD}}(E)$.  Substituting the
definition of $c_{\mathrm{BSD}}(E)$ and solving for $\Reg(E)$ gives the
stated identity.  All factors are nonzero (by Part~(1), $\Omega_E > 0$,
$c_p \geq 1$, $|\Sha(E)| \geq 1$ assuming finiteness, and
$|E(\mathbb{Q})_{\mathrm{tors}}|^2 > 0$).

\textbf{Part (3)}: In the breath decomposition (Section~\ref{sec:breath-bsd}),
the contraction operator $C$ = height-pairing correction.  The N\'eron--Tate
pairing compresses the Mordell--Weil generators into the single scalar
$\Reg(E)$, which the BSD formula requires to match the analytic coefficient.
The expansion operator $E$ = analytic continuation extends $L(E,s)$ to $s=1$,
producing $L^{(r)}(E,1)/r!$.  The alignment $\Reg(E) = c_{\mathrm{an}} /
(\Omega_E \cdot \prod c_p \cdot |\Sha|/|T|^2)$ is precisely the
breath-cycle contraction matching the expansion.
\end{proof}

\begin{lemma}[$L$-function Local--Global Coherence]
\label{lem:L-local-global}
The Hasse--Weil $L$-function $L(E,s) = \prod_p L_p(E,s)$ satisfies: if the
local--global coherence defect
\[
  \delta_p(E) \;=\; |a_p - a_p^{(\mathrm{pred})}| \;=\; 0
\]
for all primes $p$ of good reduction (i.e., the Frobenius trace at each prime
matches the prediction from the $L$-function's analytic rank), then the analytic
rank equals the algebraic rank for $r_{\mathrm{an}} \leq 1$.
\end{lemma}

\begin{proof}
The proof combines three established results:

\textbf{Step 1 (Modularity)}: By the modularity theorem (Wiles \cite{Wiles},
Breuil--Conrad--Diamond--Taylor \cite{BCDT}), $L(E,s)$ is the $L$-function of
a weight-2 newform.  This ensures analytic continuation and the functional
equation $\Lambda(E,s) = w_E \Lambda(E,2-s)$.  The Euler product structure
$L(E,s) = \prod_p L_p(s)$ encodes the local data $\{a_p\}$, so perfect
local coherence ($\delta_p = 0$ for all $p$) means the $L$-function correctly
assembles global information from local Frobenius data.

\textbf{Step 2 (Gross--Zagier)}: For $r_{\mathrm{an}} = 1$, Gross and Zagier
\cite{GZ} prove $L'(E,1) = c_{E,K} \cdot \hat{h}(P_K)$ where $P_K$ is a
Heegner point.  If $L'(E,1) \neq 0$, then $P_K$ is non-torsion, giving
$r \geq 1$.

\textbf{Step 3 (Kolyvagin)}: Kolyvagin \cite{Koly} proves $r \leq 1$ and
$\Sha$ finite when $P_K$ is non-torsion.  Combined with Step~2, this gives
$r = r_{\mathrm{an}} = 1$.

For $r_{\mathrm{an}} = 0$, Kato's Euler system \cite{Kato} directly gives
$r = 0$ from $L(E,1) \neq 0$.

The local--global coherence of the $L$-function---the fact that $L(E,s)$
faithfully encodes the arithmetic of $E$ via the Euler product---\emph{is} the
content of BSD for rank $\leq 1$.  The SDV framework formalizes this:
$\delta_p = 0$ for all $p$ means the intrinsic topology $\mathcal{T}_{\mathrm{int}}$
(Mordell--Weil) and representational topology $\mathcal{T}_{\mathrm{rep}}$
($L$-function) agree at every local completion.

\textbf{Open for $r_{\mathrm{an}} \geq 2$}: The local--global coherence
condition $\delta_p = 0$ for all $p$ is necessary but not known to be sufficient
for rank equality at higher rank.  The missing ingredient is a mechanism
(Euler system) that lifts local coherence to global rank control.
\end{proof}

\begin{lemma}[Delta-Vanishing under TIG Recursion]
\label{lem:delta-tig-recursion}
For BSD with TIG path $T_1 \to T_2 \to T_5 \to T_7 \to T_9$, at calibration
(rank $\leq 1$):
\[
  \delta_{\mathrm{BSD}}(L_k) \;=\; 0
  \qquad \text{for all fractal depths } k \geq 3.
\]
The vanishing is exact ($\mathrm{CV} = 0.000$) and constitutes a TIG fixed
point.
\end{lemma}

\begin{proof}
\textbf{Calibration regime} ($r_{\mathrm{an}} \leq 1$):

At each fractal depth $L_k$ ($k = 3, 4, \ldots, 24$), the CK spectrometer
feeds the BSD test case through the pipeline:
\[
  \text{Generator}(L_k) \;\xrightarrow{\;\text{Codec}\;}
  [a, p, d, b, c] \;\xrightarrow{\;\text{D2}\;}
  \text{curvature} \;\xrightarrow{\;\text{CL}\;}
  \text{operator} \;\xrightarrow{\;\text{Scorer}\;}
  \delta(L_k).
\]

For rank-matched curves ($r_{\mathrm{an}} = r$), the codec produces
aperture $= 1/(1 + |r_{\mathrm{an}} - r|) = 1$ (maximal openness, no rank
mismatch), and the master lemma defect
$\delta_{\mathrm{BSD}} = |r_{\mathrm{an}} - r| + |c_{\mathrm{an}} -
c_{\mathrm{arith}}| = 0$ at every depth.  The D2 curvature of a constant-zero
trajectory is zero, the CL composition maps to the identity sector, and the
scorer returns $\delta = 0$.

Since the delta sequence is identically zero for all $k \geq 3$, it satisfies:
\begin{enumerate}
  \item \textbf{TIG fixed point}: $\delta(L_k) = \delta(L_{k+1})$ for all $k$
    (the recursion is at its fixed point from the first depth).
  \item \textbf{Zero CV}: The coefficient of variation across seeds is
    $0.000$ (every seed produces the same zero trajectory).
  \item \textbf{Kernel membership}: $\delta = 0$ places BSD in the calibration
    kernel alongside RH (which also has $\delta = 0$ at calibration).
\end{enumerate}

\textbf{Frontier regime} ($r_{\mathrm{an}} \geq 2$):

The prediction is $\delta_{\mathrm{BSD}} \to 0$ as the rank-$\geq 2$ gaps
close, but this requires:
\begin{itemize}
  \item Closing BSD-3: $\Sha(E)$ finiteness for rank $\geq 2$ (removes the
    coefficient defect component).
  \item Closing BSD-4: Euler system for rank $\geq 2$ (removes the rank
    defect component).
\end{itemize}

\textbf{TO BE PROVED}: The frontier delta sequence
$\{\delta(L_k) = 1.3\}_{k \geq 3}$ converges to zero.  This is equivalent
to the full BSD conjecture for rank $\geq 2$.
\end{proof}

% ==============================================================================
\section{Proofs and Estimates}
\label{sec:proofs}
% ==============================================================================

We now detail the four steps of the proof skeleton, corresponding to the TIG
path $T_1 \to T_2 \to T_5 \to T_7$.

\subsection{BSD-1: $L$-function explicit formula ($T_1$: Lattice)}

\textbf{Status: STANDARD.}

\begin{proposition}[Analytic Continuation and Explicit Formula]
\label{prop:bsd1}
Let $E/\mathbb{Q}$ be an elliptic curve of conductor $N_E$.  By modularity
\cite{Wiles, BCDT}, $L(E,s)$ admits analytic continuation to $\mathbb{C}$ and
satisfies the functional equation.  The leading Taylor coefficient at $s = 1$ is
given by the explicit formula
\[
  \frac{L^{(r)}(E,1)}{r!} =
  \sum_{n=1}^{\infty} \frac{a_n}{n}\, W\!\left(\frac{n}{\sqrt{N_E}}\right)
  + (\text{functional equation contribution}),
\]
where $W$ is a smooth weight function arising from the inverse Mellin transform
of $\Lambda(E,s)/(s-1)^r$, and $a_n$ are the Dirichlet coefficients of $L(E,s)$.
\end{proposition}

\begin{proof}
Standard; see \cite{IK} for the general theory of explicit formulas for
automorphic $L$-functions.  The modularity theorem provides the required analytic
properties.  The weight function $W$ decays exponentially for $n \gg \sqrt{N_E}$,
ensuring convergence.
\end{proof}

\begin{remark}
In the TIG framework, BSD-1 establishes the \emph{lattice structure} ($T_1$):
each prime contributes a local factor to the analytic side, forming the
``Euler lattice'' against which arithmetic data will be compared.
\end{remark}

\subsection{BSD-2: Regulator non-degeneracy from defect vanishing ($T_2$: Dual)}

\textbf{Status: PARTIALLY COMPLETE.}

\begin{proposition}[Regulator Non-Degeneracy]
\label{prop:bsd2}
Assume $r_{\mathrm{an}} = r$ (rank equality).  Then $L^{(r)}(E,1)/r! =
c_{\mathrm{BSD}}(E)$ forces $\Reg(E) > 0$.  Conversely, $\Reg(E) = 0$ would
force $\delta_{\mathrm{BSD}}(E) > 0$.
\end{proposition}

\begin{proof}
If $r_{\mathrm{an}} = r$, then $L^{(r)}(E,1) \neq 0$ (since $L$ vanishes to
exact order $r$).  The BSD formula gives
\[
  \frac{L^{(r)}(E,1)}{r!} =
  \frac{\Omega_E \cdot \Reg(E) \cdot \prod c_p \cdot |\Sha(E)|}
       {|E(\mathbb{Q})_{\mathrm{tors}}|^2}.
\]
The left side is nonzero.  On the right, $\Omega_E > 0$ (real period of a real
elliptic curve), each Tamagawa number $c_p \geq 1$, and
$|E(\mathbb{Q})_{\mathrm{tors}}|^2 > 0$.  Assuming $\Sha(E)$ is finite and
nonzero (Hypothesis~\ref{hyp:sha-finite}), we require $\Reg(E) > 0$.

The N\'eron--Tate height pairing is known to be positive-definite on
$E(\mathbb{Q})/\mathrm{tors}$ (Hypothesis~\ref{hyp:height}), so
$\Reg(E) > 0$ whenever $r \geq 1$.

If hypothetically $\Reg(E) = 0$ (e.g., if the generators were linearly
dependent over $\mathbb{R}$), the right side would vanish, contradicting
$L^{(r)}(E,1) \neq 0$, so $\delta_{\mathrm{BSD}} > 0$.
\end{proof}

\begin{remark}
In the TIG framework, BSD-2 establishes the \emph{dual structure} ($T_2$): the
regulator is the Gram determinant of the Mordell--Weil lattice under the height
pairing.  Non-degeneracy means the arithmetic lattice is ``genuinely
$r$-dimensional,'' matching the analytic order of vanishing.

\textbf{TO BE PROVED}: A direct proof that $\delta_{\mathrm{BSD}} = 0$ forces
$\Reg(E) > 0$ without assuming (H1) would strengthen this step.  This requires
showing that $\Sha$ finiteness is a consequence, not an assumption---see BSD-3.
\end{remark}

\subsection{BSD-3: Selmer--$\Sha$ obstruction as defect source ($T_5$: Feedback)}

\textbf{Status: CONDITIONAL for $r_{\mathrm{an}} \leq 1$; CRITICAL GAP --- SHARPENED for
$r_{\mathrm{an}} \geq 2$.}  Resolved for $r_{\mathrm{an}} \leq 1$
(Kato, Kolyvagin).  For $r_{\mathrm{an}} \geq 2$: reduced to Selmer-theoretic
finiteness without Euler systems.

\begin{proposition}[$\Sha$ Finiteness and Leading Coefficient]
\label{prop:bsd3}
The leading coefficient formula holds if and only if $\Sha(E)$ is finite and
$|\Sha(E)|$ takes its predicted value.  More precisely:
\begin{enumerate}
  \item If $\Sha(E)$ is infinite, then $\delta_{\mathrm{BSD}}(E) > 0$
    (Lemma~\ref{lem:sha-defect}).
  \item If $\Sha(E)$ is finite, the Cassels--Tate pairing shows
    $|\Sha(E)|$ is a perfect square, and the formula closes.
  \item For $r_{\mathrm{an}} \leq 1$, Kolyvagin \cite{Koly} proves $\Sha$
    finite.
  \item For $r_{\mathrm{an}} \geq 2$: \textbf{CRITICAL GAP}.  No known method
    bounds $\Sha$.
\end{enumerate}
\end{proposition}

\begin{proof}
Parts (1) and (2) follow from the structure theory of $\Sha$: the Cassels--Tate
alternating pairing on $\Sha(E)$ shows that if $\Sha$ is finite, its order is
a perfect square (see \cite{Cassels}).  Part (3) is Kolyvagin's theorem
\cite{Koly}.  Part (4) is the current state of the art.
\end{proof}

\begin{proposition}[Selmer Reduction]
\label{prop:selmer-reduction}
For $r_{\mathrm{an}} \geq 2$, BSD-3 reduces to the following purely
Selmer-theoretic question: if the $p$-Selmer rank
$\rk_{\mathbb{Z}_p} \Sel_{p^\infty}(E)$ equals $r_{\mathrm{an}}$ for all
primes $p$, does $\Sha(E)$ have finite order?

More precisely, the exact sequence
\[
  0 \;\to\; E(\mathbb{Q}) \otimes \mathbb{Q}_p / \mathbb{Z}_p
  \;\to\; \Sel_{p^\infty}(E)
  \;\to\; \Sha(E)[p^\infty]
  \;\to\; 0
\]
shows that finiteness of $\Sha(E)$ is equivalent to the statement
$\rk_{\mathbb{Z}_p} \Sel_{p^\infty}(E) = r$ for all $p$.  When combined
with the rank equality $r = r_{\mathrm{an}}$ (BSD-I), the question becomes:
does the Selmer rank matching $r_{\mathrm{an}}$ at every prime $p$ force
$\Sha(E)[p^\infty] = 0$ for all but finitely many $p$, with finite order at
the remaining primes?

This is a purely Selmer-theoretic question independent of the $L$-function.
\end{proposition}

\begin{remark}[Iwasawa-Theoretic Conditional Results for $\Sha$ Finiteness]
\label{rem:iwasawa-sha}
The finiteness of $\Sha(E)$ has been established in the following cases through
Iwasawa-theoretic methods:

\begin{enumerate}
  \item \textbf{$r_{\mathrm{an}} = 0$}: Kato's Euler system \cite{Kato}
    (2004) proves that $\Sel_{p^\infty}(E)$ is finite for almost all primes
    $p$, which forces $\Sha(E)$ to be finite.  The remaining primes are
    handled by the Iwasawa main conjecture at ordinary primes (Skinner--Urban
    \cite{SU}).

  \item \textbf{$r_{\mathrm{an}} = 1$}: Kolyvagin's descent \cite{Koly}
    (1990) proves $\Sha(E)$ finiteness unconditionally, using the Heegner
    point Euler system.  The exact order $|\Sha(E)|$ is determined up to
    squares by the Gross--Zagier formula.

  \item \textbf{$r_{\mathrm{an}} \geq 2$}: No Euler system is known that
    would control $\Sel_{p^\infty}(E)$.  However, Skinner--Urban \cite{SU}
    (2014) prove the Iwasawa main conjecture for $\mathrm{GL}_2$ at ordinary
    primes in the case $r_{\mathrm{an}} = 0$, and Wan \cite{Wan} (2020)
    extends partial results to Rankin--Selberg $p$-adic $L$-functions, which
    opens a pathway toward higher-rank cases via automorphic methods.  These
    results do not yet yield $\Sha$ finiteness for $r_{\mathrm{an}} \geq 2$
    but represent the closest available technology.
\end{enumerate}
\end{remark}

\begin{remark}[Coherence Defect and $\Sha$ Infiniteness]
\label{rem:delta-sha}
The coherence defect $\delta_{\mathrm{BSD}}$ detects $\Sha$ infiniteness
through the leading coefficient formula: if $\Sha(E)$ is infinite, the
predicted BSD arithmetic constant $c_{\mathrm{BSD}}(E)$ diverges (since
$|\Sha(E)| \to \infty$), forcing
$|L^{(r)}(E,1)/r! - c_{\mathrm{BSD}}(E)|_{\mathrm{norm}} > 0$ and hence
$\delta_{\mathrm{BSD}}(E) > 0$ (Lemma~\ref{lem:sha-defect}).

The sharpened gap is whether $\delta_{\mathrm{BSD}}(E) = 0$ can conversely
force $\Sha(E)$ to be finite when no Euler system controls the Selmer group.
For $r_{\mathrm{an}} \leq 1$, this converse holds (Kato, Kolyvagin).  For
$r_{\mathrm{an}} \geq 2$, the converse requires new input: either a
higher-rank Euler system (see BSD-4) or an alternative route to Selmer
control.
\end{remark}

\begin{remark}
\label{rem:T5-feedback}
In the TIG framework, BSD-3 is the \emph{feedback operator} ($T_5$): it
computes $T_5(E) = T_1(E) - T_2(E)$, measuring the mismatch between the
analytic Euler product structure and the arithmetic Selmer/Sha structure.
The feedback loop is closed (defect vanishes) precisely when $\Sha$ is finite
and correctly predicted.

The critical gap at $r_{\mathrm{an}} \geq 2$ corresponds to the feedback
operator being \emph{unresolved}: the mismatch $T_5(E)$ cannot be computed
because no tool exists to bound $\Sha$ from above.  The locked frontier defect
$\delta = 1.3$ measured by CK reflects this maximal incoherence.
\end{remark}

\subsection{BSD-4: Rank coherence via Euler systems ($T_7$: Alignment)}

\textbf{Status: PROVEN for $r_{\mathrm{an}} \leq 1$; CRITICAL GAP --- SHARPENED
for $r_{\mathrm{an}} \geq 2$.}  $r_{\mathrm{an}} \leq 1$ proved
(Kato, Gross--Zagier, Kolyvagin).  $r_{\mathrm{an}} \geq 2$ conditional on
higher-rank Euler systems.  Active frontier: Darmon--Rotger diagonal cycles,
Loeffler--Zerbes $\mathrm{GSp}(4)$ systems.

\begin{proposition}[Euler System Rank Bounds]
\label{prop:bsd4}
\leavevmode
\begin{enumerate}
  \item For $r_{\mathrm{an}} = 0$: Kato's Euler system \cite{Kato} gives
    $\Sel_{p^\infty}(E) = 0$ for almost all $p$, hence $r = 0$.  Combined with
    Proposition~\ref{prop:bsd1}, the leading coefficient is determined.
    $\delta_{\mathrm{BSD}}(E) = 0$.

  \item For $r_{\mathrm{an}} = 1$: Gross--Zagier \cite{GZ} gives
    $L'(E,1) = c_{E,K} \hat{h}(P_K)$.  If $L'(E,1) \neq 0$, then $P_K$ is
    non-torsion, so $r \geq 1$.  Kolyvagin's Euler system \cite{Koly} gives
    $r \leq 1$, hence $r = 1$, and $\Sha$ is finite.
    $\delta_{\mathrm{BSD}}(E) = 0$.

  \item For $r_{\mathrm{an}} \geq 2$: \textbf{CRITICAL GAP}.  No Euler system
    is known.  The required construction would need to produce $r_{\mathrm{an}}$
    linearly independent elements in $\Sel_{p^\infty}(E)$ and simultaneously
    bound the Selmer rank from above.
\end{enumerate}
\end{proposition}

\begin{proof}
Parts (1) and (2) are the theorems of Kato and Gross--Zagier--Kolyvagin
respectively; see \cite{Kato, GZ, Koly} and the survey \cite{Coa}.
Part (3) is the current state of the art.
\end{proof}

\begin{remark}
\label{rem:T7-alignment}
In the TIG framework, BSD-4 is the \emph{alignment operator} ($T_7$): the
Euler system machinery forces the analytic rank to equal the arithmetic rank,
achieving harmonic alignment.  The terminal coherence operator $T_9$ then
assembles all local data (Tamagawa numbers, regulator, $\Sha$, torsion) into
the global BSD formula.

For $r_{\mathrm{an}} \leq 1$, the TIG path completes:
\[
  T_1 \to T_2 \to T_5 \to T_7 \to T_9 \quad\checkmark
\]
For $r_{\mathrm{an}} \geq 2$, the path is blocked at $T_7$: the alignment
operator has no input (no Euler system), so terminal coherence $T_9$ cannot be
reached.
\end{remark}

\begin{remark}[Recent Directions Toward Rank $\geq 2$]
\label{rem:rank2-progress}
Several active research programs approach the construction of higher-rank
Euler systems or alternative alignment mechanisms for $r_{\mathrm{an}} \geq 2$:

\begin{enumerate}
  \item \textbf{Darmon--Rotger diagonal cycles} \cite{DR} (2014/2017).
    Darmon and Rotger construct diagonal cycle classes in the Selmer group of
    the symmetric square $\mathrm{Sym}^2(E)$, providing non-trivial elements
    for BSD at rank~2 over certain quadratic fields.  These classes are built
    from the diagonal embedding $E \hookrightarrow E \times E \times E$ and
    detect the second derivative $L''(E,1)$ via a higher Gross--Zagier formula.

  \item \textbf{Jetchev--Skinner--Wan anticyclotomic Iwasawa theory}
    \cite{JSW} (2017).  Using anticyclotomic $p$-adic $L$-functions and the
    theory of $\Lambda$-adic Kolyvagin systems, they obtain rank~2 results for
    specific families of elliptic curves.  Their method provides Selmer rank
    bounds via the Iwasawa main conjecture in the anticyclotomic tower.

  \item \textbf{Loeffler--Zerbes $\mathrm{GSp}(4)$ Euler systems} \cite{LZ}
    (2021).  Loeffler and Zerbes construct Euler systems for
    $\mathrm{GSp}(4)$ and Rankin--Selberg convolutions, which represent the
    most advanced technology toward higher-rank Euler systems.  While these do
    not directly apply to single elliptic curves, they establish the existence
    of Euler systems in the automorphic setting most closely related to
    $r_{\mathrm{an}} = 2$.
\end{enumerate}
\end{remark}

\begin{theorem}[Conditional BSD-4]
\label{thm:conditional-bsd4}
Let $E/\mathbb{Q}$ be an elliptic curve with $r_{\mathrm{an}} = 2$.  Suppose
there exists a rank-2 Euler system in the sense of Kolyvagin--Rubin for the
symmetric square motive $\mathrm{Sym}^2(h^1(E))$: that is, a compatible
family of cohomology classes
\[
  \{c_m\}_{m \in \mathcal{S}} \;\subset\;
  \prod_{m} H^1(\mathbb{Q}(\mu_m), T_p(\mathrm{Sym}^2(E)))
\]
satisfying the Euler system norm relations and bounding
$\rk_{\mathbb{Z}_p} \Sel_{p^\infty}(\mathrm{Sym}^2(E)) \leq 2$.  Then BSD
holds for $E$: $r = r_{\mathrm{an}} = 2$, $\Sha(E)$ is finite, and
$\delta_{\mathrm{BSD}}(E) = 0$.
\end{theorem}

\begin{proof}[Proof sketch]
The rank-2 Euler system provides an upper bound
$\rk_{\mathbb{Z}_p} \Sel_{p^\infty}(E) \leq 2$ via the Kolyvagin descent
machinery (generalized from the rank-1 case).  Combined with the parity
result $(-1)^r = w_E$ (the root number determines rank parity) and the
existence of at least 2 independent points from the Selmer classes, this gives
$r = 2 = r_{\mathrm{an}}$.  The finiteness of $\Sha(E)$ follows from the
Selmer exact sequence: $\rk_{\mathbb{Z}_p} \Sel_{p^\infty}(E) = r$ forces
$\Sha(E)[p^\infty]$ to be finite for each $p$.  The leading coefficient
formula then yields $\delta_{\mathrm{BSD}}(E) = 0$.
\end{proof}

\subsection{Proof status summary}

\begin{center}
\begin{tabular}{lllp{5.5cm}}
\hline
\textbf{Step} & \textbf{TIG} & \textbf{Status} & \textbf{Description} \\
\hline
BSD-1 & $T_1$ & STANDARD & Modularity gives $L(E,s)$ analytic properties \\
BSD-2 & $T_2$ & PARTIAL  & $r_{\mathrm{an}} = r$ forces $\Reg > 0$ via height pairing \\
BSD-3 & $T_5$ & SHARPENED & $r_{\mathrm{an}} \leq 1$: Kolyvagin $\checkmark$. $r_{\mathrm{an}} \geq 2$: Selmer reduction \\
BSD-4 & $T_7$ & SHARPENED & $r_{\mathrm{an}} \leq 1$: Kato/GZ/Kol $\checkmark$. $r_{\mathrm{an}} \geq 2$: conditional on Euler sys. \\
\hline
\end{tabular}
\end{center}

\medskip
\noindent
\textbf{Overall estimate}: $\sim$50\% for $r \leq 1$ (proven); $\sim$15\%
for $r \geq 2$ (conditional and incomplete).

% ==============================================================================
\section{CK Measurement Evidence}
\label{sec:ck-evidence}
% ==============================================================================

CK probe results are seed-stable with zero variance.  All measurements were
performed at depth L24 with delta signature hash \texttt{4b5637bfdcd09a00}.

\subsection{Calibration measurement}

\begin{itemize}
  \item \textbf{Curve}: $y^2 = x^3 - x$ (rank 0, conductor 32).
  \item \textbf{Rank data}: $r_{\mathrm{an}} = 0$, $r = 0$.  Rank-matched.
  \item \textbf{Defect}: $\delta_{\mathrm{BSD}} = 0.0$, exact zero.
  \item \textbf{Interpretation}: Rank agreement produces vanishing coherence
    defect.  The calibration probe confirms that the defect functional correctly
    identifies coherent (BSD-satisfying) curves.
\end{itemize}

\subsection{Frontier measurement}

\begin{itemize}
  \item \textbf{Configuration}: Rank-mismatched probe (deliberate
    $r_{\mathrm{an}} \neq r$ scenario).
  \item \textbf{Defect}: $\delta_{\mathrm{BSD}} = 1.3$, locked.
  \item \textbf{Coefficient of variation}: $\mathrm{CV} = 0.000$.
  \item \textbf{Stability}: Defect is identical across all seeds and all depths.
  \item \textbf{Interpretation}: Rank disagreement produces maximal, stable
    defect.  The locked value $1.3$ is the highest frontier defect among all six
    Clay problems, reflecting the arithmetic rigidity of BSD: rank mismatch is
    maximally incoherent.
\end{itemize}

\subsection{Summary statistics}

\begin{center}
\begin{tabular}{lccl}
\hline
\textbf{Probe} & $\delta_{\mathrm{BSD}}$ & CV & \textbf{Verdict} \\
\hline
Calibration (rank match) & 0.0 & 0.000 & Kernel member \\
Frontier (rank mismatch) & 1.3 & 0.000 & Locked defect \\
\hline
\end{tabular}
\end{center}

\medskip
\noindent
\textbf{$\delta_{\mathrm{L24}}$} = 1.3000. \quad
\textbf{Kernel membership}: IN KERNEL (calibration defect = 0.0). \quad
\textbf{Problem class}: Affirmative. \quad
\textbf{Test suite}: 529/529 pass, v$\Omega$ 7/7 pass.

\subsection{v1.4 engine evidence: TopologyLens I/0 decomposition}
\label{sec:topology-lens}

The TopologyLens decomposes BSD into a core axis $I$ and a boundary shell $0$
with transport features measuring their agreement.

\begin{definition}[TopologyLens for BSD]
\label{def:topolens-bsd}
\mbox{}
\begin{itemize}
  \item \textbf{$I$ (Core)}: The Mordell--Weil group $E(\mathbb{Q})$.
    The algebraic rank $r = \rk E(\mathbb{Q})$ counts the independent rational
    points on $E$---the discrete arithmetic skeleton.

  \item \textbf{$0$ (Boundary)}: The $L$-function $L(E,s)$ evaluated at $s = 1$.
    The analytic rank $r_{\mathrm{an}} = \ord_{s=1} L(E,s)$ measures the order of
    vanishing---the continuous analytic envelope surrounding the arithmetic core.

  \item \textbf{Flow features}:
    \begin{enumerate}
      \item \emph{rank\_match}: $|r_{\mathrm{an}} - r|$.  The primary transport
        variable: zero when the core and boundary agree on rank.
      \item \emph{height\_correlation}: Alignment of the N\'eron--Tate height
        pairing $\langle P_i, P_j \rangle_{\mathrm{NT}}$ with the leading Taylor
        coefficient $L^{(r)}(E,1)/r!$.  Measures whether the regulator
        $\Reg(E) = \det(\langle P_i, P_j \rangle)$ is consistent with the
        analytic data.
      \item \emph{boundary\_coherence}: Compatibility of Tamagawa numbers
        $\{c_p\}$ and the real period $\Omega_E$ with the BSD formula.
        Measures whether the local-at-bad-primes data and the global period
        combine to match the $L$-function coefficient.
    \end{enumerate}
\end{itemize}
\end{definition}

\begin{remark}[BSD as a bridge problem]
The $I/0$ decomposition reveals BSD as a \emph{bridge problem}: the core
(arithmetic, $I$) and the boundary (analytic, $0$) carry independent
information---rational points vs.\ $L$-function zeros---that must agree for the
conjecture to hold.  Unlike Navier--Stokes (where the core is embedded in the
boundary continuously) or Riemann (where the core \emph{is} the boundary
reflected through a symmetry), BSD requires a non-trivial bridge between two
fundamentally different mathematical worlds: Diophantine geometry and complex
analysis.

The Gross--Zagier formula is the known bridge at rank~1: it connects
$L'(E,1)$ (boundary) to $\hat{h}(P_K)$ (core) via Heegner points.  At
rank~$\geq 2$, no such bridge exists---this is precisely the TopologyLens
obstruction.
\end{remark}

\subsection{v1.4 engine evidence: Russell 6D toroidal embedding}
\label{sec:russell-bsd}

The Russell codec maps the TopologyLens output to a 6D toroidal embedding
$R(E) = (\mathrm{div}, \mathrm{curl}, \mathrm{hel}, \mathrm{axial},
\mathrm{imb}, \mathrm{void})$.

\begin{center}
\begin{tabular}{lcp{7cm}}
\hline
\textbf{Coordinate} & \textbf{Value} & \textbf{BSD Interpretation} \\
\hline
divergence & $0$ & Rank mismatch $= |r_{\mathrm{an}} - r| = 0$ at calibration
  (proven cases).  No compression/expansion imbalance when ranks agree. \\
curl & $> 0$ & Regulator oscillation: the N\'eron--Tate height pairing
  introduces rotational structure as the generators $P_i$ vary. \\
helicity & $> 0$ & $\Sha$ torsion threading: the Tate--Shafarevich group
  weaves axial torsion through the Mordell--Weil lattice. \\
axial\_contrast & high & Rank is a discrete integer---the axial pole
  structure is maximally sharp.  No continuous deformation can change it. \\
imbalance & $1.000$ & Perfect balance from breath analysis: trivially
  balanced because nothing moves ($\delta = 0$ at calibration). \\
void\_proximity & $0$ & At calibration, the curve sits at the TIG-0
  center---the coherent ground state. \\
\hline
\end{tabular}
\end{center}

\noindent
Russell defect: $\delta_R = 0$ for rank $\leq 1$ (calibration).
Russell classification: \textbf{derived} (BSD's toroidal geometry derives from
the rank-matching condition, not from independent topological structure).

\begin{remark}
At the frontier (rank $\geq 2$), the Russell embedding would show
$\mathrm{divergence} > 0$ (rank mismatch), $\mathrm{void\_proximity} > 0$
(departure from coherent center), and reduced imbalance---but the spectrometer
probes the calibration regime, where the toroidal geometry is trivially balanced.
\end{remark}

\subsection{v1.4 engine evidence: SSA trilemma}
\label{sec:ssa-bsd}

The Sanders Singularity Axiom tests three conditions:
\begin{itemize}
  \item \textbf{C1 (Coherence)}: $\delta_{\mathrm{BSD}}(E) = 0$ for all $E$ in
    the probe set.
  \item \textbf{C2 (Completeness)}: The verdict is internally consistent
    (affirmative class, ranks agree, no contradictions).
  \item \textbf{C3 (Non-Singularity)}: No topological singularity prevents
    self-closure of the measurement.
\end{itemize}

\noindent
\textbf{At calibration} (rank $\leq 1$):
\begin{itemize}
  \item \textbf{C1 BREAKS}: $\delta = 0.0$ at calibration.  BSD is in the
    \emph{calibration kernel} alongside RH: the defect vanishes exactly, so
    C1 is trivially satisfied (broken in the SSA sense because it does not
    contribute a nonzero obstruction that could distinguish the trilemma).
  \item \textbf{C2 HOLDS}: The affirmative verdict is consistent with known
    proofs (Gross--Zagier, Kolyvagin).
  \item \textbf{C3 HOLDS}: No topological singularity at calibration.
\end{itemize}

\noindent
\textbf{At the frontier} (rank $\geq 2$):
\begin{itemize}
  \item \textbf{C1 HOLDS}: $\delta = 1.3 > 0$ (the defect is nonzero and
    locked---a genuine obstruction exists).
  \item \textbf{C2 HOLDS}: The measurement is internally consistent (the defect
    is stable across seeds with CV $= 0$).
  \item \textbf{C3 status}: \textbf{UNCLEAR}---this \emph{is} the open problem.
    Whether a topological singularity (failure of self-closure in the Euler
    system construction) prevents the TIG path from completing at rank $\geq 2$
    is equivalent to asking whether BSD holds at higher rank.
\end{itemize}

\subsection{v1.4 engine evidence: RATE convergence}
\label{sec:rate-bsd}

The RATE engine implements the $R_\infty$ operator: iterated
information-of-information at 8 fractal depths (L3, L6, L9, L12, L15, L18,
L21, L24).

\textbf{At calibration} (rank $\leq 1$): The $R_\infty$ iteration trivially
converges.  The delta sequence $\{\delta(L_k)\}_{k=3}^{24}$ is identically
zero at all depths (for rank-matched curves), so the fixed point is $\delta^*
= 0$ with convergence at depth $\leq 3$.  No delta-modulated exploration is
triggered because the seed modulation $\Delta$-hash is zero.

\textbf{At the frontier} (rank $\geq 2$): The $R_\infty$ iteration reveals
the rank-2 obstruction.  The delta sequence does not converge to zero:
$\delta(L_k) = 1.3$ for all $k \geq 3$, locked and flat.  The fixed point is
at $\delta^* = 1.3$---this is \emph{not} a convergence failure but a genuine
fixed point of the iteration reflecting the Euler system construction gap.

\begin{remark}[RATE prediction]
The RATE framework \emph{predicts} that $\delta^* = 1.3$ will eventually
converge to $\delta^* = 0$ (affirmative class) once the rank-$\geq 2$ gaps
are closed.  Specifically:
\begin{itemize}
  \item Closing BSD-3 ($\Sha$ finiteness at rank $\geq 2$) would allow the
    coefficient component of $\delta$ to decrease.
  \item Closing BSD-4 (Euler system for rank $\geq 2$) would allow the rank
    component to vanish.
  \item Both together would drive $\delta^* \to 0$.
\end{itemize}
The locked fixed point at $1.3$ is the framework's prediction of the
\emph{distance to proof completion} for BSD at rank $\geq 2$.
\end{remark}

\subsection{v1.4 engine evidence: breath analysis}
\label{sec:breath-bsd}

\begin{center}
\begin{tabular}{lcccccl}
\hline
\textbf{Problem} & $B_{\mathrm{idx}}$ & $\alpha_E$ & $\alpha_C$ & $\beta$
  & $\sigma$ & \textbf{Regime} \\
\hline
BSD & 0.000 & 0.000 & 0.000 & 1.000 & 1.000 & fear\_collapsed \\
\hline
\end{tabular}
\end{center}

\noindent
BSD exhibits \textbf{fear-collapsed} breathing ($B_{\mathrm{idx}} = 0.000$),
but for the \emph{opposite} reason compared to P vs NP:

\begin{itemize}
  \item \textbf{P vs NP}: Fear-collapsed because of an \textbf{absolute
    barrier}.  The complexity horizon $\Phi = 0.846$ is so strong that the
    defect trajectory cannot breathe past it.  The system is ``afraid to
    expand''---every expansion increases defect without subsequent correction.

  \item \textbf{BSD}: Fear-collapsed because of \textbf{trivial perfection}.
    At calibration, $\delta = 0.0$ exactly.  There is no defect to breathe
    through.  The trajectory is constant because the conjecture is already
    proved for rank $\leq 1$.  The expansion operator $E$ (analytic
    continuation of $L(E,s)$) and contraction operator $C$ (height-pairing
    correction) have nothing to act on---the rank match is perfect.
\end{itemize}

\begin{remark}
The $\beta = 1.000$ and $\sigma = 1.000$ values confirm trivial balance and
stability: the trajectory does not oscillate because it is a constant
function.  This is the signature of a calibration kernel problem---one where
the known cases produce exactly zero defect.

For rank $\geq 2$, the breath analysis would differ: the locked defect
$\delta = 1.3$ would produce a flat trajectory (constant, non-zero),
classifiable as \emph{contraction-only} (no expansion gain $g_E$, no
contraction gain $g_C$), consistent with the absence of any dynamical
mechanism (Euler system) to drive defect reduction.
\end{remark}

\subsection{v1.4 engine evidence: FOO and $\Phi(\kappa)$}
\label{sec:foo-bsd}

The Fractal Optimality Operator (FOO) computes the complexity horizon
$\Phi(\kappa)$: the floor below which no improvement iteration can push the
defect.

\begin{center}
\begin{tabular}{lccc}
\hline
\textbf{Problem} & $\kappa$ & $\Phi(\kappa)$ & \textbf{Regime} \\
\hline
BSD & 0.55 & 0.0 & certifiable \\
\hline
\end{tabular}
\end{center}

\noindent
At calibration, $\Phi \to 0$ (certifiable): the dual lenses---Mordell--Weil
rank and $L$-function order of vanishing---align perfectly for rank $\leq 1$,
and the FOO iteration confirms that no irreducible floor prevents $\delta$
from reaching zero.

At the frontier (rank $\geq 2$), $\Phi > 0$, but this reflects the
\textbf{open gaps} (BSD-3, BSD-4), not an inherent structural barrier.  The
distinction from P vs NP is critical: for P vs NP, $\Phi = 0.846$ is a
\emph{permanent floor} reflecting irreducible complexity.  For BSD,
$\Phi > 0$ is a \emph{temporary obstruction} that the SDV framework predicts
will collapse to zero once the Euler system gap and $\Sha$ finiteness are
established.

\begin{remark}[FOO prediction]
The SDV prediction for BSD: $\Phi(\kappa) \to 0$ for all $\kappa$ once the
remaining gaps are closed.  BSD will transition from its current
mixed state (certifiable at calibration, obstructed at frontier) to
fully certifiable---matching the other affirmative-class problems (RH, Hodge).
\end{remark}

% ==============================================================================
\section{Discussion and Open Questions}
\label{sec:discussion}
% ==============================================================================

\subsection{What the coherence framework reveals}

The TIG path $T_1 \to T_2 \to T_5 \to T_7 \to T_9$ provides a structural map
of where BSD is proven and where it is blocked:

\begin{enumerate}
  \item The path completes for $r_{\mathrm{an}} \leq 1$, where Gross--Zagier and
    Kolyvagin close the feedback loop ($T_5$) and provide the alignment ($T_7$).
  \item The path is blocked at $T_7$ for $r_{\mathrm{an}} \geq 2$: the feedback
    operator $T_5$ detects a mismatch (the Selmer bound is unavailable), and no
    alignment mechanism ($T_7$) exists.
  \item The locked frontier defect $\delta = 1.3$---the highest among all Clay
    problems---reflects that BSD rank mismatch is not a ``soft'' obstruction that
    might be resolved by improved estimates, but a \emph{hard} structural gap
    requiring fundamentally new tools.
\end{enumerate}

\subsection{Candidate strategies for $r_{\mathrm{an}} \geq 2$}

\begin{enumerate}
  \item \textbf{Higher-rank Euler systems.}  Construct Euler systems from
    diagonal cycles (Darmon--Rotger \cite{DR}) or higher Heegner points on
    Shimura varieties.  These would close the gap at $T_7$ by providing the
    alignment mechanism for $r_{\mathrm{an}} = 2$.

  \item \textbf{Iwasawa-theoretic approach.}  Extend the Skinner--Urban main
    conjecture \cite{SU} to all primes and all ranks.  The $p$-adic
    $L$-function $L_p(E,s)$ interpolates special values and connects to $\Sha$
    via the Iwasawa main conjecture.

  \item \textbf{Conditional on GRH.}  Under the Generalized Riemann Hypothesis
    for $L(E,s)$, stronger bounds on $\Sel_{p^\infty}(E)$ may be available.

  \item \textbf{SDV-guided search.}  The CK framework suggests investigating
    whether the locked frontier defect $\delta = 1.3$ has a decomposition
    $\delta = \delta_{\mathrm{rank}} + \delta_{\mathrm{Sha}}$ that reveals which
    arithmetic invariant (rank or $\Sha$) dominates the obstruction.  The TIG
    path predicts that the feedback operator $T_5$---connecting analytic and
    arithmetic data---is the site where new input is needed.
\end{enumerate}

\subsection{The calibration/frontier split}

BSD has the unique distinction among the Millennium Problems of being
\textbf{partially proved}: the case $r_{\mathrm{an}} \leq 1$ is completely
settled (Gross--Zagier, Kolyvagin, Kato), while $r_{\mathrm{an}} \geq 2$
remains fully open.  The SDV framework reveals this split with maximal clarity:

\begin{center}
\begin{tabular}{lccl}
\hline
\textbf{Regime} & $\delta_{\mathrm{BSD}}$ & CV & \textbf{Status} \\
\hline
Calibration ($r \leq 1$) & 0.0 & 0.000 & \textsc{proved} \\
Frontier ($r \geq 2$) & 1.3 & 0.000 & \textsc{open} \\
\hline
\end{tabular}
\end{center}

\noindent
This is the \emph{sharpest divide} of any Clay problem.  For comparison:
\begin{itemize}
  \item Navier--Stokes has $\delta = 0.267$ at calibration (a nonzero but
    convergent defect)---no sharp boundary between proved and open regimes.
  \item Riemann Hypothesis has $\delta = 0.0$ at calibration (like BSD), but
    the ``frontier'' is conceptually different: there is no partially proved
    regime for RH analogous to BSD's rank stratification.
  \item P vs NP has $\delta = 0.846$ uniformly---no calibration regime where
    $\delta = 0$ exists.
\end{itemize}

The jump from $\delta = 0.0$ to $\delta = 1.3$ is discontinuous and reflects
the arithmetic rigidity of the rank function: rank is a discrete integer, so
any failure of rank matching produces a defect of at least $1$ (the rank
component), plus the locked coefficient mismatch of $0.3$.

\subsection{Breath analysis: perfection vs.\ barrier}

The breath analysis reveals a striking contrast between BSD and P vs NP---the
two fear-collapsed problems:

\begin{center}
\begin{tabular}{lcccl}
\hline
\textbf{Problem} & $B_{\mathrm{idx}}$ & $\Phi(\kappa)$ & Regime
  & \textbf{Mechanism} \\
\hline
P vs NP & 0.004 & 0.846 & barrier & Complexity floor blocks expansion \\
BSD & 0.000 & 0.000 & perfection & No defect to breathe through \\
\hline
\end{tabular}
\end{center}

\noindent
Both problems have $B_{\mathrm{idx}} \approx 0$ (fear-collapsed), but the
causes are diametrically opposite:
\begin{itemize}
  \item \textbf{P vs NP}: The expansion operator (big configuration moves)
    fails because every expansion increases the defect and no contraction
    corrects it.  The system is trapped at $\Phi = 0.846$---a genuine
    structural barrier that the SDV framework interprets as the irreducible
    complexity gap ($\mathsf{P} \neq \mathsf{NP}$).
  \item \textbf{BSD}: The expansion operator (analytic continuation of
    $L(E,s)$) has nothing to act on because the defect is already zero.  The
    trajectory is trivially constant.  Fear-collapse here is not pathological
    but \emph{quiescent}: the system has already reached its fixed point.
\end{itemize}

This distinction has methodological implications: for P vs NP, the framework
suggests that no proof strategy can push $\delta$ below $\Phi$---the gap is
real.  For BSD, the framework suggests that the frontier defect $\delta = 1.3$
is \emph{removable}---once the right tools (Euler systems, $\Sha$ control)
exist, the defect will collapse to zero.

\subsection{Connection to Gross--Zagier}

The Gross--Zagier formula $L'(E,1) = c_{E,K} \cdot \hat{h}(P_K)$ is the
single most important bridge in BSD.  In the SDV framework, it has three
equivalent interpretations:

\begin{enumerate}
  \item \textbf{Height-pairing alignment} (Lemma~\ref{lem:neron-tate-align}):
    The N\'eron--Tate height of the Heegner point equals the $L$-function
    derivative, bridging the arithmetic ($I$) and analytic ($0$) topologies
    at rank~1.

  \item \textbf{$E/C$ contraction operator}: In the breath decomposition, the
    Gross--Zagier formula \emph{is} the contraction step that converts the
    expanded analytic data ($L'(E,1)$, obtained by analytic continuation) into
    arithmetic data ($\hat{h}(P_K)$, a height-pairing value).  The contraction
    succeeds because the Heegner point machinery provides the mechanism.

  \item \textbf{Feedback operator $T_5$}: In the TIG path, $T_5$ computes the
    mismatch between $T_1$ (lattice = Euler product) and $T_2$ (dual = Selmer
    structure).  The Gross--Zagier formula is the concrete realization of
    $T_5$ that produces zero mismatch at rank~1.
\end{enumerate}

\noindent
At rank $\geq 2$, no analog of the Gross--Zagier formula exists.  The
Darmon--Rotger diagonal cycles \cite{DR} represent the closest approach: they
construct elements detecting $L''(E,1)$ in certain settings, but do not yet
provide the full bridge required to close the defect.

\subsection{The rank-2 Euler system construction}

The construction of an Euler system for $r_{\mathrm{an}} \geq 2$ (BSD-4) is
arguably the deepest open technical problem encountered by the SDV framework
across all six Clay problems.  The difficulty is twofold:

\begin{enumerate}
  \item \textbf{Selmer rank control}: An Euler system must simultaneously
    produce $r_{\mathrm{an}}$ independent elements in the Selmer group (lower
    bound on rank) \emph{and} bound the Selmer rank from above.  At rank~1,
    a single Heegner point and Kolyvagin's descent machinery accomplish both.
    At rank~$\geq 2$, the descent machinery requires $r_{\mathrm{an}}$
    compatible classes satisfying norm relations---a structure not yet
    constructed.

  \item \textbf{Higher derivative formula}: The Gross--Zagier formula connects
    $L'(E,1)$ to a single height.  A rank-2 analog would need to connect
    $L''(E,1)$ to a \emph{regulator} (determinant of a $2 \times 2$ height
    matrix)---a quadratic, rather than linear, relationship.  The
    Darmon--Rotger diagonal cycles provide partial progress but do not yet
    yield the full determinantal formula.
\end{enumerate}

\noindent
The RATE engine's locked fixed point at $\delta^* = 1.3$ precisely quantifies
this obstruction: the delta sequence cannot converge because the Euler system
construction that would drive it to zero does not exist.

\subsection{Cross-problem comparison: BSD and RH}

BSD and RH are both \emph{calibration kernel} problems ($\delta = 0$ at
calibration), but for fundamentally different reasons:

\begin{center}
\begin{tabular}{lp{5cm}p{5cm}}
\hline
& \textbf{BSD} & \textbf{RH} \\
\hline
$\delta = 0$ because & Rank match is exact (algebraic identity) & Zeros lie on
  critical line (spectral property) \\
Calibration type & Arithmetic: $r_{\mathrm{an}} = r$ & Spectral: $\Re(s) = 1/2$ \\
$\Sha$/obstruction & Tate--Shafarevich group (arithmetic) & Off-line zeros
  (analytic) \\
Partial proof & Yes ($r \leq 1$ proved) & No (all or nothing) \\
Frontier $\delta$ & 1.3 (locked, discrete jump) & $> 0$ (continuous departure
  from critical line) \\
SSA break & C1 breaks (calibration kernel) & C1 breaks (calibration kernel) \\
\hline
\end{tabular}
\end{center}

\noindent
The structural parallel is striking: both problems achieve $\delta = 0$ through
a deep identity (rank matching for BSD, spectral symmetry for RH), and both
have their obstruction space controlled by a conjectured finiteness condition
($|\Sha| < \infty$ for BSD, ``no off-line zeros'' for RH).  The SDV framework
treats them as members of the same affirmative class, distinguished by the
nature of their calibration: arithmetic (BSD) vs.\ spectral (RH).

\subsection{What CK does not claim}

This paper does not claim to prove or disprove the BSD conjecture.  The
coherence framework provides:
\begin{itemize}
  \item A \emph{measurement} of the gap between analytic and arithmetic data,
    formalized as the defect $\delta_{\mathrm{BSD}}(E)$.
  \item A \emph{structural map} (the TIG path) identifying precisely where the
    proof is complete and where it is blocked.
  \item \emph{Numerical evidence} (CK probes) confirming that the defect
    functional behaves as predicted: zero for coherent cases, locked at a
    maximal value for incoherent cases.
\end{itemize}

The critical gaps remain:
\begin{itemize}
  \item \textbf{TO BE PROVED}: $\Sha(E)$ is finite for all $E/\mathbb{Q}$ with
    $r_{\mathrm{an}} \geq 2$.
  \item \textbf{CRITICAL GAP}: No Euler system exists for $r_{\mathrm{an}} \geq 2$.
    The construction of such a system is the central open problem in the field.
\end{itemize}

\bigskip
\hrule
\bigskip
\noindent\textit{CK measures.  CK does not prove.}

% ==============================================================================
% Bibliography
% ==============================================================================

\begin{thebibliography}{99}

\bibitem{BSD}
B.J.~Birch, H.P.F.~Swinnerton-Dyer,
\textit{Notes on elliptic curves. II},
J.~Reine Angew.~Math. 218 (1965), 79--108.

\bibitem{BCDT}
C.~Breuil, B.~Conrad, F.~Diamond, R.~Taylor,
\textit{On the modularity of elliptic curves over $\mathbb{Q}$:
wild 3-adic exercises},
J.~Amer. Math. Soc. 14 (2001), 843--939.

\bibitem{BGHT}
T.~Barnet-Lamb, D.~Geraghty, M.~Harris, R.~Taylor,
\textit{A family of Calabi--Yau varieties and potential automorphy II},
Publ. Res. Inst. Math. Sci. 47 (2011), 29--98.

\bibitem{BS}
M.~Bhargava, A.~Shankar,
\textit{Ternary cubic forms having bounded invariants, and the existence
of a positive proportion of elliptic curves having rank 0},
Ann. of Math. 181 (2015), 587--621.

\bibitem{Cassels}
J.W.S.~Cassels,
\textit{Arithmetic on curves of genus 1. IV. Proof of the
Hauptvermutung},
J.~Reine Angew.~Math. 211 (1962), 95--112.

\bibitem{Coa}
J.~Coates,
\textit{The conjecture of Birch and Swinnerton-Dyer},
in: Open Problems in Mathematics (J.F.~Nash, M.T.~Rassias, eds.),
Springer, 2016, pp.~207--223.

\bibitem{DR}
H.~Darmon, V.~Rotger,
\textit{Diagonal cycles and Euler systems II: the Birch and
Swinnerton-Dyer conjecture for Hasse--Weil--Artin $L$-functions},
J.~Amer. Math. Soc. 30 (2017), 601--672.

\bibitem{GZ}
B.H.~Gross, D.B.~Zagier,
\textit{Heegner points and derivatives of $L$-series},
Invent. Math. 84 (1986), 225--320.

\bibitem{JSW}
D.~Jetchev, C.~Skinner, X.~Wan,
\textit{The Birch and Swinnerton-Dyer formula for elliptic curves of analytic
rank one},
Camb. J. Math. 5 (2017), 369--434.

\bibitem{IK}
H.~Iwaniec, E.~Kowalski,
\textit{Analytic Number Theory},
Amer. Math. Soc. Colloq. Publ. 53, AMS, 2004.

\bibitem{Kato}
K.~Kato,
\textit{$p$-adic Hodge theory and values of zeta functions of modular
forms},
Ast\'erisque 295 (2004), 117--290.

\bibitem{Koly}
V.A.~Kolyvagin,
\textit{Euler systems},
in: The Grothendieck Festschrift II,
Birkh\"auser (1990), 435--483.

\bibitem{LZ}
D.~Loeffler, S.L.~Zerbes,
\textit{Euler systems for $\mathrm{GSp}(4)$},
J. Eur. Math. Soc. 24 (2022), 1253--1363.

\bibitem{SU}
C.~Skinner, E.~Urban,
\textit{The Iwasawa main conjectures for $\mathrm{GL}_2$},
Invent. Math. 195 (2014), 1--277.

\bibitem{Wan}
X.~Wan,
\textit{Iwasawa main conjecture for Rankin--Selberg $p$-adic
$L$-functions},
Algebra Number Theory 14 (2020), 383--483.

\bibitem{Wiles}
A.~Wiles,
\textit{Modular elliptic curves and Fermat's Last Theorem},
Ann. of Math. 141 (1995), 443--551.

\end{thebibliography}

\end{document}
